% Korean Noether paper 24 — fresh producer translation, unchecked.
% Initial translation authority: ED0005 lines 13658--14146.
% Current scope-rebase authority: ED0011 / NOETH-DE-AUTH-v047-20260808.
% Authority SHA-256: 1A44F967B29972E8F99E5C323A479162AD82A23FC457395915A4BB9DDF51AD41.
% Normalized source span: 98,379 bytes; SHA-256 04ACA866CB66D6F504EC1016C41FFF5B4D6A16B8F97DD28983AEB8C4B1E12D22.
% Translation only; independent source, Korean-language, render, and certification checks remain external.

\providecommand{\frakS}{\mathfrak S}
\providecommand{\fraka}{\mathfrak a}
\providecommand{\frakb}{\mathfrak b}
\providecommand{\frakc}{\mathfrak c}
\providecommand{\frakg}{\mathfrak g}
\providecommand{\frakG}{\mathfrak G}
\providecommand{\frakh}{\mathfrak h}
\providecommand{\frakm}{\mathfrak m}
\providecommand{\frakM}{\mathfrak M}
\providecommand{\frakn}{\mathfrak n}
\providecommand{\frako}{\mathfrak o}
\providecommand{\frakp}{\mathfrak p}
\providecommand{\frakq}{\mathfrak q}
\providecommand{\frakr}{\mathfrak r}
\providecommand{\frakt}{\mathfrak t}
\providecommand{\frK}{\mathfrak K}
\providecommand{\barfrakm}{\overline{\mathfrak m}}

\section*{24. 소거이론과 일반 아이디얼 이론}
\begin{center}
Math. Ann. 90 (1923), 229--261쪽\\[0.9em]
{\Large\bfseries 소거이론과 일반 아이디얼 이론}\\[1.0em]
괴팅겐의 에미 뇌터\\[0.6em]
\rule{4em}{0.4pt}
\end{center}

이하에서는 내가 헨첼의 서술에 부여한 산술적 형태의 소거이론\footnote{Kurt Hentzelt, \emph{Zur Theorie der Polynomideale und Resultanten}. E. Noether 편, Math. Ann. 88 (1922), 53--79쪽. 이하 H.-N.으로 인용한다.} --- 이는 다항식영역의 아이디얼 이론이라고 할 수 있다 --- 을 일반 아이디얼 이론\footnote{E. Noether, \emph{Idealtheorie in Ringbereichen}, Math. Ann. 83 (1921), 24--66쪽. 이하 \emph{아이디얼 이론}으로 인용한다.} 안에 배치하고, 동시에 영점에 관한 부분을 새로 정초한다. 두 이론의 기본개념과 체론의 기본개념\footnote{E. Steinitz, \emph{Algebraische Theorie der Körper}, J. f. M. 137 (1910), 167--309쪽. 이하 Steinitz로 인용한다.}은 § 1에 짧게 모아 두었으므로 여기서 참조할 수 있다.

이 배치와 새로운 정초는 네 문제를 중심으로 한다. 차원개념의 산술적 정식화, 영점이론, 노름과 초등인자의 분해정리와 일반 아이디얼 이론의 분해정리 사이의 관계, 그리고 노름과 초등인자형에 의한 소아이디얼과 으뜸아이디얼의 특징짓기가 그것이다.

\emph{차원개념}의 산술적 정식화(§ 4)는 소아이디얼의 사슬로 주어진다. 이는 곡면 위에 곡선이 있고 곡선 위에 점이 있다는 사실의 산술적 등가물이다. 소거이론의 \emph{매개변수 정의}와 일치함이 증명되는 이 산술적 정식화는 조금만 고치면 임의의 환영역으로 옮길 수 있다. 나머지 문제들의 어떤 부분들을 함께 옮기면 거기서 아이디얼의 구조를 정확히 통찰할 수 있으며, 이는 다른 곳에서 보이겠다.

H.-N.에서는 통상적인 방식대로 주어진 아이디얼의 \emph{영점} 문제를 연속소거로 돌아가 직접 다룬다. 이 때문에 검산의 성격을 완전히 피하지는 못한다. 이와 대조되는 것은 한 부정원 다항식에 늘 쓰는 방법이다. 주어진 다항식을 \emph{소함수들의 거듭제곱들의 곱}(으뜸함수들의 곱)으로 나타내고, 각 소함수에 대하여 잉여류체와 \emph{동형인} 영점체를 구성한다. 소함수의 차수는 이 체의 차수를 주고, 그 영점이 소함수의 영점과 일치하는 으뜸함수의 차수는 이 으뜸함수에 따른 잉여류환의 차수를 준다. 각 소함수에 관하여 갈루아체로 넘어가면 일차인수분해를 얻는다. 이로써 원래 다항식의 영점, 일차인수분해, 중복도 문제가 해결된다. 이와 완전히 나란하게 여기서는(§ 3) 한 소아이디얼에 따른 잉여류계를 분수몫을 만들어 잉여류체로 확대하고, 그와 동형인 영점체를 구성한다. 이 영점체는 \(x_{i+1},\ldots,x_n\) 같은 부정원들을 첨가하여 생성한 부분체를 포함한다. 노름의 차수는 이 부분체에 관한 차수를 주며, 동시에 갈루아체로 넘어가면 초등인자형과 따라서 노름의 일차인수분해를 얻는다. 한편 대응하는 으뜸아이디얼의 영점들은 같다. 노름이 소아이디얼의 초등인자형의 한 거듭제곱이 되므로(§ 2) 분해도 함께 주어진다. 노름의 차수는 으뜸아이디얼에 따른 잉여류환이 부정원들에서 유도된 부분체에 관하여 갖는 차수를 준다.

이에 따라 임의의 아이디얼의 영점 문제는 \emph{노름과 초등인자의 분해정리와 일반 아이디얼 이론의 분해정리 사이의 관계}(§ 5)로 환원된다. 아이디얼의 영점들이 그 개별 대응 소아이디얼들의 영점들로 이루어진다는 사실에는 노름을 \emph{극대 으뜸인자}들로 분해하는 것이 대응한다. 이 인자들은 각 대응 소아이디얼의 초등인자형의 거듭제곱과 같다.\footnote{소거이론과 일반 아이디얼 이론의 이 평행성은 크로네커의 소거이론에서는 성립하지 않는다. H.-N., 주 4)를 보라.} 이로써 노름의 일차인수분해가 일반적으로 이루어진다. H.-N.에서는 중복도를 기초아이디얼들과 그 성분들로 해석한다. 노름의 한 극대 으뜸인자의 차수는 \(x_{i+1},\ldots,x_n\)을 첨가했을 때 \((i-1)\)단계 기초아이디얼을 \(i\)단계 기초아이디얼의 한 성분으로 나눈 선형독립 잉여류들의 수와 같다. 이제 \((i-1)\)단계 기초아이디얼은 \((n-i)\)차원보다 높은 모든 소아이디얼들에 의해 유일하게 정의되는 고립성분과 같음이 증명된다. \(i\)단계 기초아이디얼의 한 성분은 노름인자에 대응하는 소아이디얼과 더 높은 차원의 모든 소아이디얼들에 의해 정의되는 고립성분과 같다. 그리고 대응하는 노름인자의 차수는 \(x_{i+1},\ldots,x_n\)을 첨가했을 때, 이 성분의 잉여류들 가운데 \((i-1)\)단계 기초아이디얼의 류들만으로 이루어진 부분환의 차수와 같다.

\emph{소아이디얼과 진으뜸아이디얼의 특징짓기}(§ 6)는 잉여류체를 고찰하면 직접 따라온다. 원래 계수영역이 완전체이면 소아이디얼에는 노름과 초등인자형으로 소함수들이, 진으뜸아이디얼에는 진으뜸함수들이 대응하며, 역도 성립한다. 계수영역이 비완전체이면 필요조건 또는 충분조건만 줄 수 있다. 예들이 보여 주듯 한 소아이디얼의 노름이 진으뜸함수가 될 수도 있고, 한 진으뜸아이디얼의 초등인자형이 소함수가 될 수도 있다. 더 정확히 말하면 표수 \(p\)에서 소아이디얼의 노름은 그 초등인자형의 \(p^g\)제곱(\(g\geqq0\))이 된다. 반면 초등인자형이 소함수인 진으뜸아이디얼의 노름은 임의의 거듭제곱이 될 수 있으며, 이 역시 예들이 보여 준다. 끝으로 \emph{절대소아이디얼}(§ 7), 곧 계수영역을 대수적으로 닫힌 체로 확대해도 소아이디얼로 남는 아이디얼들을 고찰한다. 한 유한 대수적 수체의 계수를 갖는 절대소아이디얼의 경우, 위 특징짓기에서 A. Ostrowski가 절대소함수에 대하여 처음 세운 정리의 일반화가 나온다. 유한개의 예외를 제외하면 수체의 모든 소아이디얼을 법으로 해도 소아이디얼이라는 성질이 유지된다.

이 마지막 문제와 대수적 수체의 아이디얼 이론 사이의 유사성도 지적해 두자. 그러한 아이디얼의 초등인자형은 그 아이디얼로 나누어지는 가장 작은 양의 유리정수라고 보아야 한다(주 10). 따라서 소아이디얼의 경우에는 언제나 소수가 된다. 역으로 한 아이디얼의 노름이 소수인 순간 그 아이디얼은 소아이디얼이다. 증명도 완전히 평행하다. 그러므로 위에서 \emph{말한 예들}은 비완전체에서 고차 소아이디얼과 분기아이디얼의 유사물이 나타나는 반면, 완전체에서는 1차 소아이디얼만 있고 분기아이디얼은 없음을 보인다.

\subsection*{§ 1. 소거이론·일반 아이디얼 이론·체론의 기본개념}

\textbf{1. 영역.} \(\bar{\frakS}\)를 추상적으로 정의된 체 \(P\)의 계수를 갖는 \(y_1,\ldots,y_n\)의 모든 다항식 정역이라 하자. \(y\)들에 부정원 \(u_{\mu\nu}\)를 계수로 갖는 변환을 가한다.\footnote{§ 3 등을 고려하여 이 변환은 H.-N.의 것보다 조금 더 일반적이며, 따라서 그곳의 모든 결과는 더욱 그대로 유지된다.}
\srcnumdisplay{(1)}{y_1=u_{11}x_1+\cdots+u_{1n}x_n;\ \ldots;\ y_n=u_{n1}x_1+\cdots+u_{nn}x_n}
그 역변환은
\srcnumdisplay{(1')}{x_1=t_{11}y_1+\cdots+t_{n1}y_n;\ \ldots;\ x_n=t_{1n}y_1+\cdots+t_{nn}y_n}
이다. 서로 유리적으로 나타낼 수 있으므로 같은 뜻인 \(u_{\mu\nu}\), 또는 \(t_{\mu\nu}\)를 \(P\)에 첨가하여 계수영역 \(P(u)=P(t)\)를 얻는다. \(\frakS\)를 \(P(u)\)의 계수를 갖는 \(x_1,\ldots,x_n\)의 모든 다항식 정역이라 하자. 영역 \(\frakS\)와 \(\bar{\frakS}\)가 소거이론의 바탕이고, \(\frakS\)의 아이디얼 \(\fraka,\frakb,\ldots\)와 \(\bar{\frakS}\)의 아이디얼 \(\bar{\fraka},\bar{\frakb},\ldots\)를 고찰한다. 임의의 환에서 아이디얼은 통상적으로 \(\alpha,\beta\)와 함께 \(\alpha-\beta\)를 포함하고, \(\alpha\)와 함께 \(\lambda\alpha\)도 포함하는 것으로 정의한다. 여기서 \(\lambda\)는 임의의 환원소이다. \(\frako\)는 언제나 모든 원소로 이루어진 단위아이디얼을 뜻한다.

\textbf{2. 변환된 아이디얼.} \(\frakS\)의 아이디얼 \(\frakm\)이 \(\bar{\frakS}\)의 한 아이디얼 \(\barfrakm\)에서 (1)에 따라 \(u_{\mu\nu}\), 또는 \(t_{\mu\nu}\)를 첨가하여 생길 때 이를 변환된 아이디얼이라 한다. \(\bar f(y)\), 또는 \(\bar g(y)\)가 \(\barfrakm\)의 모든 다항식을 돌 때 \(\frakm\)은 모든 선형결합
\[
 f(x)=\sum U_\lambda(u)\bar f_\lambda(y)=\sum U_\lambda(u)f_\lambda(x)
\]
또는
\[
 g(x)=\sum T_\lambda(t)\bar g_\lambda(y)=\sum T_\lambda(t)g_\lambda(x)
\]
로 이루어진다. 특히 모든 \(f_\lambda(x)=\bar f_\lambda(y)\), 또는 같은 뜻으로 모든 \(g_\lambda(x)=\bar g_\lambda(y)\)가 들어 있다. \(f(x)\)와 \(g(x)\)는 \(P(u)=P(t)\)의 양만큼을 제외하고 정해지므로, 일반성을 잃지 않고 \(U_\lambda,T_\lambda\)를 각각 \(u,t\)의 거듭제곱곱이라고 가정할 수 있다. \(U,T\)를 임의의 다른 거듭제곱곱으로 바꾸어도 \(f(x),g(x)\)는 다시 \(\frakm\)의 다항식으로 간다. 특히 (1)의 열들, 또는 (1')의 행들을 맞바꾸어도 그러하다. 따라서 \(\frakm\)은 임의의 한 다항식과 함께 \(x\)들을 맞바꾸어 얻는 모든 다항식을 포함한다. 다만 \(u\) 또는 \(t\)에 의존하는 계수들에서도 대응하는 치환을 해야 한다. 이하에 나오는 모든 아이디얼은 반대말을 명시하지 않는 한 변환된 것이라고 가정한다. 변환되지 않은 아이디얼도 (1)과 \(x_i=v_{i1}z_1+\cdots+v_{in}z_n\)을 합성하면 \(P(u,v)\)의 계수를 갖는 \(z\)의 다항식영역 안의 변환된 아이디얼로 간다. 반면 변환된 아이디얼에서는 이 합성은 \(u_{\mu\nu}\)를 \(u,v\)의 쌍선형 결합들로 바꾸는 것일 뿐이다.

\textbf{3. 표기.} \(f^{(i)},g^{(i)},\ldots,a^{(i)},b^{(i)},\ldots\)는 시종일관 \(\frakS\)의 다항식 가운데 \(x_1,\ldots,x_{i-1}\)을 포함하지 않는 것들을 뜻한다.

\textbf{4. 기초아이디얼.} 각 아이디얼 \(\frakm\)에는 다음 규정에 따라 \(n\)개의 기초아이디얼 \(\frakg_0,\frakg_1,\ldots,\frakg_{n-1}\)이 대응한다. \((i-1)\)단계 기초아이디얼 \(\frakg_{i-1}\)은 \(b^{(i)}G(x)\equiv0(\frakm)\)이 되도록 하는 0 아닌 다항식 \(b^{(i)}\)가 존재하는 모든 \(G(x)\), 그리고 오직 그러한 다항식들만을 포함한다. 일반적으로 \(b^{(i)}\)는 \(G(x)\)에 따라 달라진다. 아이디얼 기저의 존재에서 모든 \(G(x)\)에 공통으로 쓸 수 있는 \(B^{(i)}\)가 적어도 하나 존재함도 따라온다. 따라서
\begin{equation}
 B^{(i)}\frakg_{i-1}\equiv0(\frakm),\qquad B^{(i)}\ne0 . \tag{2}
\end{equation}
변환된 아이디얼의 기초아이디얼도 변환된 아이디얼이다(H.-N. 정리 VI). \(i\)단계 기초아이디얼 \(\frakg_i\)는 \(x_{i+1},\ldots,x_n\)을 \(P(u)\)에 첨가했을 때 \(\frakm\)이 옮겨가는 아이디얼로도 정의된다. 이때 \(x_{i+1},\ldots,x_n\)에 관하여 \emph{다항식인} 것들로 제한해도 일반성을 잃지 않는다. \(\frakg_0\)은 \emph{언제나} 단위아이디얼 \(\frako\)와 같다.

\textbf{5. 아이디얼의 모듈표현, 초등인자와 개별노름.} 각 다항식 \(f(x)\)을 \(x_1,\ldots,x_{i-1}\)의 거듭제곱곱 \(\xi\)에 관한 선형형식으로 보고 다항식 \(a^{(i)}\)를 계수로 삼으면, 아이디얼 \(\frakm\)은 \(a^{(i)}\)들의 영역에 관한 선형형식들의 가군 \(\frakM_{i-1}\)으로 간다. 곧 \(\frakM_{i-1}\)은 임의의 두 선형형식과 함께 그 차를 포함하고, 한 선형형식 \(l(\xi)\)와 함께 임의의 \(a^{(i)}\)에 대한 \(a^{(i)}l(\xi)\)도 포함한다. 그러한 가군 \(\frakM\)의 기초모듈은 \(b^{(i)}g(\xi)=0(\frakM)\), \(b^{(i)}\ne0\)이 되도록 하는 모든 선형형식 \(g(\xi)\)의 집합이다. 따라서 모듈표현에서 기초아이디얼 \(\frakg_{i-1}\)은 \(\frakM_{i-1}\)의 기초모듈 \(\frakG_{i-1}\)으로 간다. \(\frakM_{i-1}\)과 \(\frakG_{i-1}\)은 무한히 많은 부정원 \(\xi\)에 의존하지만, 각 선형형식에는 그 가운데 유한개만 나타난다. \(\frakM_{i-1}\)은 \(x_{i+1},\ldots,x_n\)을 \(P(u)\)에 첨가했을 때 \emph{단위원과 다른 초등인자를 유한개만 갖는} 특징적인 성질을 가진다. 곧 이 첨가 뒤 \(\frakG_{i-1}\)과 \(\frakM_{i-1}\)은 다음 기저표현을 갖는다. 여기서 \(\zeta\)들은 새 부정원이며, \(\xi\)들과 \(x_i\)에 관하여 다항식인 가역 선형변환으로 이어진다. \(\zeta_i=r_i(\xi)\), \(\xi_k=s_k(\zeta)\)에는 매번 이 부정원들 가운데 유한개만 나타난다.
\begin{align}
\frakG_{i-1}&=(\zeta_0,\zeta_1,\ldots,\zeta_\sigma,\zeta_{\sigma+1},\ldots,\zeta_{\sigma+\nu},\ldots), \notag\\
\frakM_{i-1}&=(E^{(i)}\zeta_0,E_1^{(i)}\zeta_1,\ldots,E_\sigma^{(i)}\zeta_\sigma,
\zeta_{\sigma+1},\ldots,\zeta_{\sigma+\nu},\ldots), \tag{3}
\end{align}
여기서 각 \(E_\mu^{(i)}\)는 그 앞의 것을 나눈다(H.-N. (24), (28)).

\emph{최고 초등인자} \(E^{(i)}\)는 (2)에 나타나는 모든 \(B^{(i)}\)의, 다항식 의미에서의 최대공약인자로도 정의된다. \(x_{i+1},\ldots,x_n\)에 관하여 다항식이고 원시적이라고 잡을 수 있으며, 그러면 \(x_i\)에 관하여 정칙이다. 곧 \(x\)들에 관한 차수가 \(\lambda\)이면 0 아닌 계수를 갖는 항 \(x_i^\lambda\)를 포함한다.

(3)에 나타나는 모든 초등인자의 곱 \(R^{(i)}\)을 \(\frakG_{i-1}\)에서 \(\frakM_{i-1}\)로 가는 노름, 또는 아이디얼 \(\frakm\)의 \emph{개별노름}이라 한다. \(x_{i+1},\ldots,x_n\)을 첨가하면 \(\frakm\)이 \(\frakg_i\)로 간다는 점을 고려하여 기호로 쓰면
\begin{equation*}
 R^{(i)}=E^{(i)}E_1^{(i)}\cdots E_\sigma^{(i)}
 =N(\frakG_{i-1}\mid\frakM_{i-1})=N(\frakg_{i-1}\mid\frakg_i).
\end{equation*}
(3)이 보여 주듯 \(x_{i+1},\ldots,x_n\)을 첨가하면 \(R^{(i)}\)은 \(\frakG_{i-1}\)에서 \(\frakM_{i-1}\)로 가는 전이치환의 행렬식과 같다. \(R^{(i)}\)의 차수는 \(P(u,x_{i+1},\ldots,x_n)\)에 관하여 선형독립인 \(\frakG_{i-1}\)에 따른 \(\frakM_{i-1}\)의 잉여류들, 따라서 \(\frakg_{i-1}\)에 따른 \(\frakg_i\)의 잉여류들의 수를 준다. \(R^{(i)}\)은 \(x_{i+1},\ldots,x_n\)에 관하여 다항식이고 원시적이라고 잡을 수 있으며, 그러면 \(x_i\)에 관하여 정칙이다. 또한 \(R^{(i)}\)은 다음 성질을 갖고 언제나 존재하는 랭크 \(\varrho\)의 가군 \(\frakM^*_{i-1}\)의 모든 \(\varrho\)행 행렬식들의 다항식 의미의 최대공약인자로 계산할 수 있다. \(\frakM^*_{i-1}\)의 기초모듈을 \(\frakG^*_{i-1}\)이라 할 때, 잉여류계 \(\frakG^*_{i-1}\mid\frakM^*_{i-1}\)은 \(\frakG_{i-1}\mid\frakM_{i-1}\)과 동형이다(H.-N. 정리 VII과 § 5).

\textbf{6. 아이디얼의 초등인자형과 노름(종결식형).} 모든 최고 초등인자의 곱 \(E_\frakm=E^{(1)}E^{(2)}\cdots E^{(n)}\)을 \(\frakm\)의 \emph{초등인자형}, 모든 개별노름의 곱 \(R_\frakm=R^{(1)}R^{(2)}\cdots R^{(n)}\)을 \emph{노름(종결식형)}\footnote{H.-N.에서는 종결식형이라는 명칭만 쓴다. 노름이라는 명칭은 산술적 성질에 대응한다.}이라 한다. 초등인자형과 노름은 아이디얼로 나누어지고, 노름은 초등인자형으로 나누어지며, 초등인자형의 어떤 거듭제곱은 노름으로 나누어진다. 더 일반적으로
\begin{align}
 E^{(i)}\frakg_{i-1}&\equiv0(\frakg_i),
& E^{(n)}\cdots E^{(i)}\frakg_{i-1}&\equiv0(\frakm),\notag\\
 R^{(i)}\frakg_{i-1}&\equiv0(\frakg_i),
& R^{(n)}\cdots R^{(i)}\frakg_{i-1}&\equiv0(\frakm), \tag{4}
\end{align}
이 성립한다(H.-N. 정리 VIII). 더 나아가 \(\frakm\)이 \(\frakn\)으로 나누어지고 노름 \(R_\frakm,R_\frakn\)이 일치하면 \emph{아이디얼 \(\frakm,\frakn\)도 일치한다}(H.-N. 정리 IX). \(\frakm\)의 한 인자에 대하여 \(R^{(1)},R^{(2)},\ldots\)가 일치하면 기초아이디얼 \(\frakg_1,\frakg_2,\ldots\)도 일치하고 언제나 \(\frakg_0=\frako\)임을 고려하면, 이에 맞추어 다음을 얻는다. \(\frakn\)이 \(\frakm\)의 \emph{진인자}이면 \(\frakn\)의 개별노름 가운데 \(\frakm\)의 대응하는 개별노름과 처음으로 다른 것은 \emph{진인자}이다.\footnote{반면 \(R_\frakn\)의 뒤쪽 인자들은 \(R_\frakm\)의 대응 인자들의 배수가 될 수 있다. 예컨대 \(\frakm\)이 소아이디얼이고 \(\frakn\)이 단위아이디얼이 아닐 때에는 언제나 그렇다.} \emph{\(\frakm\)이 0 아닌 다항식 \(G^{(i)}\)를 포함하면} 정의상 \(\frakg_0,\frakg_1,\ldots,\frakg_{i-1}\)은 단위아이디얼과 같다. 따라서 개별노름이 선형독립 잉여류들의 수를 정한다는 5의 성질에 의해 \(R^{(1)},\ldots,R^{(i-1)}\)은 \emph{단위원}과 같고, 그러므로 \(E^{(1)},\ldots,E^{(i-1)}\)도 \emph{단위원}과 같다.

\textbf{7. 개별노름과 초등인자의 분해, 기초아이디얼의 성분들.} \(P(u)\)에 관하여 기약인 다항식을 소함수, 소함수의 거듭제곱을 으뜸함수라 한다. 동시에 소함수가 아닌 으뜸함수, 곧 1보다 높은 거듭제곱을 진으뜸함수라 한다. 한 다항식의 극대 으뜸인자들로의 분해란 각 인자가 으뜸함수이지만 임의의 두 인자의 곱은 더 이상 으뜸이 아닌 분해이다.\footnote{소함수와 으뜸함수의 정의는 11과 정확히 나란히 아이디얼이라는 말만 다항식으로 바꾸어도 된다. 극대 으뜸인자들은 12의 분해에서 극대 으뜸성분들에 대응한다.} 개별노름 \(R^{(i)}=N(\frakg_{i-1}\mid\frakg_i)\)을 극대 으뜸인자 \(Q_\lambda^{(i)}\)들로 분해하는 데에는 \(\frakg_i\)를 최소공배수
\[
\frakg_i=[\frakr_1,\frakr_2,\ldots,\frakr_s]
\]
로 나타내는 것이 대응하며, \(Q_\lambda^{(i)}=N(\frakg_{i-1}\mid\frakr_\lambda)\)이다.
이때 \(\frakr_\lambda\)는 \(\frakg_{i-1}\)을 \((i-1)\)단계 기초아이디얼로 가지며 그 \(i\)단계 기초아이디얼과 일치한다. \(\frakr_\lambda\)들을 \emph{기초아이디얼의 성분들}이라 한다. \(\frakr_\lambda\)의 최고 초등인자는 \(Q_\lambda^{(i)}\)와 \(E^{(i)}\)의 다항식 의미의 최대공약인자, 따라서 \(E^{(i)}\)의 한 극대 으뜸인자와 같다. \(\frakr_\lambda\)들은 위에서 말한 기초아이디얼들에 관한 거동과, 개별노름 또는 최고 초등인자가 각각 \(R^{(i)}\) 또는 \(E^{(i)}\)의 극대 으뜸인자가 되어야 한다는 요구에 의해 유일하게 결정된다(H.-N. 정리 X).

\textbf{8. 차원개념.} \(R^{(i)}\)이 단위원과 다른 첫 개별노름일 때 아이디얼 \(\frakm\)의 최고차원을 \(n-i\)라 한다. 6에 따라 같은 뜻으로 \(\frakg_i\)가 단위아이디얼과 다른 첫 기초아이디얼일 때이다. 단위아이디얼 \(\frako\)의 차원은 \(-1\)이라 한다. 단위원과 다른 개별노름 \(R^{(i)}\)이 하나뿐이어서 \(\frakg_i=\frakg_{i+1}=\cdots=\frakm\)이면 최고차원을 단순히 \(\frakm\)의 차원이라고도 한다. \(\frakm\)이 0 아닌 다항식 \(G^{(i)}\)을 포함하면 6의 끝에 있는 주에서 알 수 있듯 그 최고차원은 기껏해야 \(n-i\)이다. \(\frakm\)의 최고차원이 \(n-i\)이고 \(\frakn\)이 \(\frakm\)의 인자이면 \(\frakn\)의 최고차원도 기껏해야 \(n-i\)이다.

\textbf{9. 노름(종결식형)과 영점.} 아이디얼 \(\frakm\)의 영점이란 \(P(u;x_{i+1},\ldots,x_n)\)의 어떤 적당한 대수적 확대체에 속하는 \(x_1,\ldots,x_i\)의 값의 계들(\(i=1,2,\ldots,n\)) 가운데, \(x_{i+1},\ldots,x_n\)을 계수영역에 첨가했을 때 \(\frakm\)의 모든 다항식이 사라지는 것들을 뜻한다. 노름(종결식형)의 각 인자 \(R^{(i)}\)에는 이 첨가 뒤 그 차수만큼의 \(\frakm\)의 영점이 대응한다. \(R^{(i)}\)을 \(u_{\mu\nu}\)와 다른 부정원 \(v\)의 쌍선형 결합 \(w_{\mu\nu}\)로 쓰면 다음 명시적 분해를 갖는다.
\begin{equation}
 R^{(i)}=\prod\{z-(v_1\bar{x}_{1\nu}+\cdots+v_{i-1}\bar{x}_{i-1,\nu}+v_i\bar{x}_{i,\nu})\}^{\lambda_\nu}. \tag{5}
\end{equation}
여기서 \(\bar{x}_{1\nu},\ldots,\bar{x}_{i,\nu}\)는 서로 대응하는 영점들의 한 계이다(H.-N. 정리 XIII. § 3에서 새 증명을 줄 것이다). 따라서 최고차원 \(n-i\)인 아이디얼의 영점들 가운데 \((n-i)\)개 매개변수에 의존하는 것들은 있지만, 그보다 더 많은 매개변수에 의존하는 것은 없다. 이로써 8에서 준 차원개념이 통상적인 정식화와 일치함을 보였다.

\textbf{10. 일반 아이디얼 이론의 환영역.} 바탕 영역은 \emph{유한사슬 정리}가 성립하는 가환환이라 하자. 곧 아이디얼의 뜻에서 각 \(\fraka_i\)가 바로 앞의 아이디얼의 진인자인 모든 사슬 \(\fraka_1,\fraka_2,\ldots,\fraka_k,\ldots\)은 유한 단계에서 끝난다. 이 요구는 모든 아이디얼이 아이디얼 기저를 갖는다는 요구와 동등하므로(\emph{아이디얼 이론}, 정리 I), 특히 다항식영역에서 만족한다. 이어지는 11, 12, 13에서는 이 일반 환영역을 바탕으로 삼는다.

\textbf{11. 소아이디얼과 으뜸아이디얼.} 한 곱이 \(\frakp\)로 나누어지면 적어도 한 인자가 \(\frakp\)로 나누어지는 아이디얼 \(\frakp\)를 소아이디얼이라 한다. 한 곱이 \(\frakq\)로 나누어지면 한 인자 또는 각 인자의 어떤 거듭제곱이 \(\frakq\)로 나누어지는 아이디얼 \(\frakq\)를 으뜸아이디얼, 또는 으뜸이라 한다. 기호로 쓰면
\[
 \quad a\not\equiv0(\frakp),\quad b\not\equiv0(\frakp)\quad\text{이면}\quad ab\not\equiv0(\frakp);
\]
\[
 \quad a\not\equiv0(\frakq),\quad b^x\not\equiv0(\frakq)\ \text{가 모든 거듭제곱 \(x\)에 성립하면}\quad ab\not\equiv0(\frakq).
\]
각 으뜸아이디얼 \(\frakq\)에는 유일한 대응 소아이디얼 \(\frakp\)가 있다. 이는 \(\frakq\)의 인자이고 그 어떤 거듭제곱은 \(\frakq\)로 나누어진다. 곧 \(\frakq\equiv0(\frakp)\), \(\frakp^\rho\equiv0(\frakq)\)이다.
이때 \(\frakp\)는 그 어떤 거듭제곱이 \(\frakq\)로 나누어지는 모든 환원소로 이루어진다. 정의가 보여 주듯 모든 소아이디얼은 동시에 으뜸아이디얼이다. 동시에 소아이디얼이 아닌 으뜸아이디얼, 따라서 \(\rho>1\)인 것을 진으뜸아이디얼이라 한다.

\textbf{12. 분해정리.} 아이디얼을 최소공배수로 나타낸 표현 \(\fraka=[\frakq_1,\ldots,\frakq_s]\)에서 어느 \(\frakq\)도 나머지 것들의 최소공배수, 곧 그 보충부분으로 나누어지지 않을 때 이를 최단표현이라 한다. 각 \(\frakq\)가 으뜸이지만 임의의 두 \(\frakq\)의 최소공배수가 으뜸이 아닐 때 극대 으뜸성분들에 의한 표현이라 한다. 모든 아이디얼은 \emph{유한개의 극대 으뜸성분에 의한 최단표현}을 갖는다. 서로 다른 두 그러한 표현에서는 \emph{성분들의 수와, 모두 서로 다른 대응 소아이디얼들이 일치한다}. 이렇게 유일하게 결정되는 소아이디얼들을 그 아이디얼의 소아이디얼 또는 대응 소아이디얼이라 한다(\emph{아이디얼 이론}, 정리 IX).

\textbf{13. 분해의 고립성분.} \(\frakr=[\frakq_{i_1},\ldots,\frakq_{i_\lambda}]\)가 \(\fraka\)의 극대 으뜸성분들에 의한 적어도 하나의 최단표현에 나타나는 \(\frakq_{i_j}\)들로 이루어지고, 이들의 대응 소아이디얼 가운데 어느 것도 \(\fraka\)의 나머지 소아이디얼들 가운데 하나로 나누어지지 않는다면 \(\frakr\)을 \(\fraka\)의 분해의 고립성분이라 한다. 분해의 \emph{고립성분들은 그 소아이디얼들에 의해 유일하게 결정된다}. 특히 고립된 극대 으뜸성분들은 유일하게 결정된다(\emph{아이디얼 이론}, 정리 XIII).

\textbf{14. 표수, 완전체와 비완전체, 제1종 확대.} \(\Omega\)에 포함되고 단위원에서 유도되는 소체가 유리수체의 유형이면 \(\Omega\)의 표수는 0이고, 이 소체가 소수 \(p\)에 따른 잉여류계의 유형이면 표수는 \(p\)라고 한다. \(\Omega\)의 계수를 갖는 모든 소함수가 어떤 적당한 확대체에서 서로 다른 일차인자들로 분해되면 \(\Omega\)를 완전체라 하고, 그렇지 않으면 비완전체라 한다. 표수 0인 모든 체는 완전체이다. 표수 \(p\)인 체가 완전체인 것은 각 원소와 함께 그 \(p\)제곱근도 포함할 때 그리고 오직 그때이다. 비완전체의 한 소함수는 \(x^{p^f}\)에 관한 \(r\)차 다항식이고, 어떤 적당한 확대체에서 서로 다른 \(r\)개 일차인자의 \(p^f\)제곱들로 분해된다. \(f\)를 지수라 한다. 어떤 적당한 확대체에서 서로 다른 일차인자들로 분해되어 지수 \(f\)가 0인 소함수를 제1종이라 한다. 모든 원소가 제1종, 곧 제1종 소함수의 영점인 확대를 제1종 확대라 한다. 제1종 원소를 첨가하여 생기는 모든 확대는 제1종이고, 완전체에는 제1종 확대만 있다(Steinitz, § 11과 § 13).

\textbf{15. 유한확대의 축약차수와 지수.} \(\frK\)가 비완전체 \(\Omega\)의 유한확대이면, \(\frK\) 안에는 \(\Omega\)에 관하여 제1종이고 \(\frK\)의 모든 제1종 원소, 그리고 오직 그 원소들만으로 이루어진 차수 \(r\)의 부분체 \(\frK_0\)가 있다. \(\frK\)의 모든 원소의 \(p^f\)제곱은 \(\frK_0\)에 속하고, 정확히 차수 \(rp^f\)인 원소들도 \(\frK\) 안에 있다. \(r\)을 축약차수, \(f\)를 이 유한확대의 지수라 한다(Steinitz, § 14, 1과 § 14, 5).

\subsection*{§ 2. 소아이디얼과 으뜸아이디얼의 초등인자형과 노름. 소아이디얼의 인자}

이제부터는 시종일관 § 1, 1에서 정의한 다항식영역을 다룬다. 먼저 소아이디얼과 으뜸아이디얼도 변환된 아이디얼로 제한하여 다룰 수 있음을 보여야 한다.

\textbf{보조정리 I.} \emph{\(\bar{\frakS}\)의 한 소아이디얼 \(\bar{\frakp}\)을 변환하여 얻은 아이디얼 \(\frakp\)는 소아이디얼이고, \(\bar{\frakS}\)의 한 으뜸아이디얼 \(\bar{\frakq}\)을 변환하여 얻은 \(\frakq\)는 으뜸아이디얼이다. 달리 말하면 소아이디얼 또는 으뜸아이디얼이라는 성질은 \(u_{\mu\nu}\)를 \(P\)에 첨가해도 유지된다.}

\(\fraka\not\equiv0(\frakp)\), \(\frakb\not\equiv0(\frakp)\)라고 하자. 여기서 \(\fraka,\frakb\)는 반드시 변환된 아이디얼일 필요는 없다. 예컨대 \(\fraka,\frakb\)의 다항식 \(a(x),b(x)\)를 각각 \(a(x)\not\equiv0(\frakp)\), \(b(x)\not\equiv0(\frakp)\)이도록 택하자. (1)을 (1')로 역변환하면, 일반성을 잃지 않고 \(T\)들을 \(t_{\mu\nu}\)의 거듭제곱곱으로 두어
\[
 a(x)=\sum T_j \bar a_j(y),\qquad b(x)=\sum T_j \bar b_j(y)
\]
를 얻는다. 적어도 하나의 \(\bar a_j(y)\)와 하나의 \(\bar b_j(y)\)는 각각 \(\bar{\frakp}\)로 나누어지지 않아야 한다. \(T\)들을 어떤 사전식 순서로 정렬했을 때 \(\bar a_r(y)\not\equiv0(\bar{\frakp})\), \(\bar b_s(y)\not\equiv0(\bar{\frakp})\)가 각각 최고항이라 하자. 그러면 \(\bar a_r(y)\bar b_s(y)\not\equiv0(\bar{\frakp})\)이고, 따라서 \(a(x)b(x)\not\equiv0(\frakp)\), 곧 \(\fraka\frakb\not\equiv0(\frakp)\)이다. 이로써 \(\frakp\)가 소아이디얼임을 보였다. 이 증명은 소아이디얼에 따른 잉여류계가 영인자 없는 환이라는 사실, 그리고 그 성질이 부정원을 첨가해도 유지된다는 사실에 바탕을 둔다.

이에 대응하여 \(\fraka\not\equiv0(\frakq)\), 그리고 모든 거듭제곱 \(z\)에 대하여 \(\frakb^z\not\equiv0(\frakq)\)라고 하자. 아이디얼 기저가 존재하므로 \(\fraka\)의 어떤 \(a(x)\)와 \(\frakb\)의 어떤 \(b(x)\)를 \(a(x)\not\equiv0(\frakq)\), 그리고 모든 \(z\)에 대하여 \(b(x)^z\not\equiv0(\frakq)\)이도록 택할 수 있다. 그렇지 않다고 가정하면 \(\frakb\)의 어떤 거듭제곱이 \(\frakq\)로 나누어질 것이기 때문이다.

\(a(x)b(x)=c(x)\)라 두고, \(\fraka^*,\frakb^*,\frakc^*\)를 각각 \(a(x),b(x),c(x)\)의 계수 \(\bar a_j(y),\bar b_j(y),\bar c_j(y)\)에서 유도되는 아이디얼이라 하자. 모든 거듭제곱 \(z\)에 대하여 \(\frakb^{*z}\not\equiv0(\bar{\frakq})\)이다. 그렇지 않으면 \(b(x)^z\)가 \(\frakq\)로 나누어지기 때문이다. \(\fraka^*\not\equiv0(\bar{\frakq})\)이므로 모든 거듭제곱 \(\lambda\)에 대하여 \(\fraka^*\frakb^{*\lambda}\not\equiv0(\bar{\frakq})\)이어야 한다. 그런데 데데킨트--메르텐스 정리\footnote{H.-N., 63쪽과 주 13)을 보라.}에 따라 \(\fraka^*\frakb^{*h}=\frakb^{*h-1}\frakc^*\)이 되게 하는 지수 \(h\)가 존재한다. 따라서 \(\frakc^*\not\equiv0(\bar{\frakq})\)이어야 한다. 이는 \(c(x)\not\equiv0(\frakq)\), 따라서 \(\fraka\frakb\not\equiv0(\frakq)\)을 준다. 이로써 \(\frakq\)가 으뜸임을 보였다.

한 아이디얼을 최소공배수로 나타낼 때에도 변환된 아이디얼들로 제한할 수 있다.

\textbf{보조정리 II.} \emph{\(\bar{\frakS}\)에서 \(\barfrakm=[\bar{\frakc}_1,\ldots,\bar{\frakc}_\alpha]\)라는 표현이 있으면 \(\frakS\)에서 \(\frakm=[\frakc_1,\ldots,\frakc_\alpha]\)라는 표현을 얻는다. 여기서 \(\frakc_i\)들은 \(\bar{\frakc}_i\)들을 변환하여 얻은 아이디얼들이다. 따라서 극대 으뜸성분들에 의한 모든 최단표현(§ 1, 12)에서도 이 성분들을 변환된 아이디얼들이라고 가정할 수 있다.}

\(\frakm\)의 모든 다항식 \(f(x)\)은 필요하면 \(u_{\mu\nu}\)의 적당한 거듭제곱을 곱한 뒤 \(f(x)=\sum U_i\bar f_i(y)\) 꼴이고, 각 \(\bar f_i(y)\)는 \(\barfrakm\)의 다항식으로서 가정상 모든 \(\bar{\frakc}_i\)로 나누어지므로, \(\frakm\)은 \([\frakc_1,\ldots,\frakc_\alpha]\)로 나누어진다. 역으로 \(f(x)\)가 각 \(\frakc_i\)로 나누어지면 위 꼴이고, 따라서 \(\frakm\)으로 나누어진다. 그러므로 \(\bar{\frakS}\)의 극대 으뜸성분들에 의한 분해에서 출발하면 \(\frakS\)에서 분해 \(\frakm=[\frakq_1,\ldots,\frakq_\alpha]\)를 얻는다. 여기서 \(\frakq\)들은 보조정리 I에 따라 으뜸아이디얼들의 변환이므로 으뜸이다. 또 \(\bar{\frakS}\)의 서로 다른 대응 소아이디얼 \(\bar{\frakp}\)들에는 \(\frakS\)의 서로 다른 \(\frakp\)들이 대응하므로 이들은 극대 으뜸성분들이다.

이제 소아이디얼과 으뜸아이디얼을 초등인자형과 노름으로 처음 특징짓는다. 아직 소인 것과 으뜸인 것을 완전히 분리하지는 않는다. 대수적 수체의 아이디얼 이론과 정확히 나란히 다음이 성립한다.\footnote{대수적 수체의 아이디얼에서 최고 초등인자는 그 아이디얼로 나누어지는 가장 작은 양의 유리정수라고 보아야 한다. 따라서 특히 소아이디얼의 경우에는 그 소아이디얼로 나누어지는 소수이고, 으뜸아이디얼(소아이디얼의 거듭제곱)의 경우에는 소수의 거듭제곱이다. 실제로 모든 아이디얼 \(\fraka^*\)는 동시에 체기저 \(\omega_1,\ldots,\omega_n\)의 원소들에 관한 선형형식들의 가군 \(A^*\)이고, \((\omega_1,\ldots,\omega_n)\)은 \(A^*\)의 기초모듈이다. 따라서 \(\fraka^*\)로 나누어지는 가장 작은 양의 유리정수가 가군 \(A^*\)의 최고 초등인자이고, \(\fraka^*\)의 노름은 \(A^*\)의 모든 초등인자의 곱이다.}

\textbf{정리 I.} \emph{한 소아이디얼의 초등인자형은 소함수이고, 그 노름은 이 소함수의 한 거듭제곱이다. 으뜸아이디얼의 경우 초등인자형과 노름은 으뜸함수이고, 더 정확히는 같은 소함수의 거듭제곱들이다. 이 소함수 자체는 대응 소아이디얼의 초등인자형이다. 소아이디얼과 으뜸아이디얼에서는 최고차원이 곧 차원이고, 으뜸아이디얼과 그 대응 소아이디얼의 차원은 일치한다.}

단위아이디얼에서는 초등인자형과 노름이 단위원과 같으므로 정리 I이 분명히 성립한다. 이제 소아이디얼 \(\frakp\)의 최고차원이 \(n-i\)라고 하자. § 1, 6의 (4)에 따라
\[
 E^{(n)}\cdots E^{(i+1)}\frakg_i\equiv0(\frakp),\quad\text{그러나}\quad
 E^{(n)}\cdots E^{(i+1)}\not\equiv0(\frakp)
\]
이다. 그렇지 않으면 § 1, 8에 따라 \(\frakp\)의 최고차원은 기껏해야 \(n-i-1\)일 것이기 때문이다. 따라서 \(\frakg_i\equiv0(\frakp)\)이고, \(\frakp\)는 그 \(i\)단계 기초아이디얼, 그러므로 \(\frakg_{i+1},\frakg_{i+2},\ldots\)와도 같다. 따라서 \(R^{(i+1)}=N(\frakg_i\mid\frakg_{i+1}),R^{(i+2)},\ldots\)은 단위원과 같다. 그 인자인 \(E^{(i+1)},E^{(i+2)},\ldots\)도 단위원과 같다. 그러므로 초등인자형은 \(E^{(i)}\)과 같고, 이것은 한 소함수 \(P^{(i)}\)와 같아야 한다. 실제로 § 1, 5에 따라 \(\frakg_{i-1}\)이 단위아이디얼임을 고려하면 \(E^{(i)}\)은 \(\frakp\)로 나누어지는 모든 \(B^{(i)}\)의 다항식 의미의 최대공약인자이다. \(E^{(i)}=C_1^{(i)}C_2^{(i)}\equiv0(\frakp)\)이면 \(E^{(i)}\)은 그 인자 \(C_1^{(i)}\), 또는 \(C_2^{(i)}\) 가운데 하나로 나누어져야 한다. 또 \(E^{(i)}\)의 어떤 거듭제곱은 \(R^{(i)}\)으로 나누어지므로 \(R^{(i)}\)은 이 소함수 \(P^{(i)}\)의 한 거듭제곱이며, 첫 거듭제곱일 수도 있다. 동시에 \(R^{(i)}\)은 \(\frakp\)의 노름이다. 단위원과 다른 개별노름 \(R^{(i)}\)이 하나뿐이므로 끝으로 \(\frakp\)의 최고차원은 곧 차원이다.

이에 대응하여 으뜸아이디얼 \(\frakq\)의 최고차원이 \(n-i\)라고 하자. 그러면 위와 같이
\[
 E^{(n)}\cdots E^{(i+1)}\frakg_i\equiv0(\frakq),\quad
 (E^{(n)}\cdots E^{(i+1)})^z\not\equiv0(\frakq)
\]
가 모든 거듭제곱 \(z\)에 대하여 성립한다. 따라서 위와 같이 \(\frakq=\frakg_i\)이고, 그 초등인자형은 \(E^{(i)}\), 노름은 \(R^{(i)}\)이며, 최고차원은 곧 차원이다. 이 차원은 대응 소아이디얼의 차원과 같다. 실제로 \(\frakp^\rho\equiv0(\frakq)\), \(\frakq\equiv0(\frakp)\)에서 \(\frakp\)의 차원 \(n-i\)은 기껏해야 \(\frakq\)의 차원이고, \(\frakq\)의 차원은 기껏해야 \(\frakp^\rho\)의 차원임이 따른다. \(P^{(i)}\)이 \(\frakp\)의 초등인자형일 때 \((P^{(i)})^\rho\equiv0(\frakp^\rho)\)이므로 마지막 최고차원은 기껏해야 \(n-i\)이다. 따라서 \(n-i\)은 \(\frakp\)와 \(\frakq\)의 차원이다. 더 나아가 \((P^{(i)})^\rho\equiv0(\frakq)\)에서 \(E^{(i)}\)은 \(\frakq\)의 모든 \(B^{(i)}\)의 최대공약인자로서 \(P^{(i)}\)의 한 거듭제곱이며, 첫 거듭제곱일 수도 있다. 따라서 \(R^{(i)}\)에도 같은 사실이 성립한다. 이로써 정리 I의 모든 부분을 증명하였다.

차원개념에 의한 소아이디얼의 특징은 다음 정리가 준다.

\textbf{정리 II.} \emph{소아이디얼은 같은 최고차원을 갖는 진인자를 갖지 않는다. 역으로 그러한 진인자를 갖지 않는 모든 아이디얼은 소아이디얼이다.}

증명은 다음 보조정리에 바탕을 둔다.

\textbf{보조정리 III.} \emph{한 체 \(\Omega\)의 계수를 갖는 다항식들의 소아이디얼 \(\frakp\)가 \(\Omega\)에 관하여 선형독립인 잉여류를 유한개만 가지면, \(\frakp\)는 단위아이디얼과 다른 진인자를 갖지 않는다. 달리 말하면 \(\frakp\)에 따른 잉여류계는 체이다.}

가정은 다음을 뜻한다. 임의의 \(k+1\)개 잉여류 사이에는 \(\Omega\)의 계수, 더 정확히는 \(\Omega\)의 모든 원소가 나타내는 잉여류체 \((\Omega)\)의 계수를 갖는 선형종속관계가 있지만, \((\Omega)\)에 관하여 선형독립인 \(k\)개의 잉여류로 이루어진 계가 적어도 하나 있다. 여기서 모든 잉여류를 임의의 선형독립인 \(k\)개 계로 선형표현할 수 있음이 곧바로 따른다. 이제 \(\fraka\)가 \(\frakp\)의 진인자이고, \(a(z)\not\equiv0(\frakp)\)인 다항식 \(a(z)\)가 \(\fraka\)에 속한다고 하자. \(c_1f_1(z)+\cdots+c_kf_k(z)\not\equiv0(\frakp)\)이면
\[
 c_1a(z)f_1(z)+\cdots+c_ka(z)f_k(z)\not\equiv0(\frakp)
\]
도 성립한다. 곧 \(f_1,\ldots,f_k\)의 잉여류들뿐 아니라 \(af_1,\ldots,af_k\)의 잉여류들도 \(k\)개의 선형독립인 계를 이루며, 특히 단위원류를 이들로 선형표현할 수 있다. 따라서 \(\fraka\)는 단위아이디얼과 같다. 달리 말하면 잉여류계는 체이다. 실제로 \(a(z)\not\equiv0(\frakp)\)마다 \(1\equiv a(z)f(z)(\frakp)\)이 되는 \(f(z)\)가 존재한다.

정리 II의 첫 부분을 증명하려면 차원 \(n-i\)인 모든 소아이디얼, 더 일반적으로 최고차원 \(n-i\)인 모든 아이디얼이 \(P(u;x_{i+1},\ldots,x_n)\)에 관하여 선형독립인 잉여류를 유한개만 가진다는 점만 주목하면 된다. \(\frakg_{i-1}\)이 여기서 단위아이디얼이므로 그 수는 (§ 1, 5) \(R^{(i)}\)의 차수가 주는 수이다. 따라서 \(\frakp\)의 모든 진인자 \(\fraka\)는 \(x_{i+1},\ldots,x_n\)을 \(P(u)\)에 첨가하면 단위아이디얼로 간다. 달리 말하면 \(\fraka\)의 \(i\)단계 기초아이디얼은 단위아이디얼이고, 따라서 \(\fraka\)의 최고차원은 기껏해야 \(n-i-1\)이다. 변환되지 않은 아이디얼 \(\fraka\)도 인자로 허용하려면 § 1, 2에 따라 새로운 변환된 영역으로 넘어가기만 하면 된다.

역으로 \(\frakp\)가 같은 최고차원 \(n-i\)을 갖는 진인자를 갖지 않는다고 하자. \(\fraka\not\equiv0(\frakp)\), \(\frakb\not\equiv0(\frakp)\)이고 \(\fraka,\frakb\)는 다시 변환된 아이디얼이라고 하자. 아이디얼 의미의 최대공약인자 \((\fraka,\frakp)\)와 \((\frakb,\frakp)\)는 \(\frakp\)의 진인자이므로 가정에 따라
\[
 G^{(i+1)}\equiv0(\fraka,\frakp),\qquad
 H^{(i+1)}\equiv0(\frakb,\frakp),\qquad
 G^{(i+1)}\not\equiv0,\qquad
 H^{(i+1)}\not\equiv0
\]
을 얻는다. 따라서
\[
 G^{(i+1)}H^{(i+1)}\equiv0(\fraka\frakb,\frakp),\qquad
 G^{(i+1)}H^{(i+1)}\not\equiv0.
\]
그러므로 \(\fraka\frakb\equiv0(\frakp)\)이면 \(\frakp\)의 최고차원이 \(n-i\)보다 작아질 것이므로 반드시 \(\fraka\frakb\not\equiv0(\frakp)\)이다. \(\fraka,\frakb\)가 변환된 아이디얼이므로 \(\bar{\frakp}\), 그리고 보조정리 I에 따라 \(\frakp\)도 소아이디얼이다. 이로써 정리 II를 증명하였다.

\subsection*{§ 3. 소아이디얼과 으뜸아이디얼의 영점}

먼저 소아이디얼에 관하여 영점이론(§ 1, 9)을 직접 정초하는 길을 열어 주는 것은 보조정리 III을 확장한 다음 보조정리이다. 이는 임의의 환의 소아이디얼에 대하여 성립한다.

\textbf{보조정리 IV.} \emph{단위아이디얼과 다른 임의의 소아이디얼에 따른 잉여류계는 분수몫을 만들어 한 체, 곧 그 소아이디얼의 잉여류체로 확대할 수 있다.}

임의의 아이디얼에 따른 잉여류계는 환을 이룬다. 합과 차와 곱이 다시 그 계에 속하고, 결합법칙·교환법칙·분배법칙은 잉여류로 넘어가도 유지되기 때문이다. 특별히 소아이디얼이면 이 환에는 영인자가 없다. 소아이디얼의 정의(§ 1, 11)가 사라지지 않는 잉여류들의 곱도 사라지지 않는다고 말하기 때문이다. 소아이디얼이 단위아이디얼과 다르면 이 환에는 영원소와 다른 원소가 적어도 하나 있으므로, 분수몫을 만드는 것, 곧 원소쌍들을 첨가하는 것으로 체까지 확대할 수 있다.\footnote{Steinitz, § 3. 원래 정역에 영과 다른 원소가 \emph{하나} 있다는 가정으로 충분하다. 그러면 분수몫을 만들 때 단위원의 존재가 보장되어 원래 영역이 그 체의 부분영역이 되기 때문이다. Steinitz는 정역에 단위원이 존재한다고 가정한다.} 이로써 잉여류체가 정의된다.

이하에서는 다음 \textbf{표기}를 일관되게 쓴다. 잉여류는 그 류의 임의의 한 원소를 괄호에 넣어 \((x),\ldots,(a(x)),\ldots\)처럼 나타내고, 잉여류들로 이루어진 체도 괄호로 나타낸다. 특히 \((P)\), \((P(u))\), \((P(u,x_{i+1},\ldots,x_n))\)는 각각 \(P\), \(P(u)\), \(P(u,x_{i+1},\ldots,x_n)\)의 원소들로 나타낼 수 있는 잉여류들 모두와 오직 그 잉여류들로 이루어진 잉여류체를 뜻한다.

또 다음을 미리 말해 둔다. 한 체 \(\mR\)이 바탕체 \(\Omega\)에 관하여 대수적 랭크(초월차수) \(k\)라는 것은, \(\mR\)이 \(\Omega\)에 관하여 대수적으로 독립인 \(k\)개 원소, 곧 초월원소들로 이루어진 계를 적어도 하나 포함하지만 \(\mR\)의 임의의 \(k+1\)개 원소는 \(\Omega\)에 관하여 대수적으로 종속이라는 뜻이다. \(\Omega\)에 관하여 대수적으로 독립인 \(z_1,\ldots,z_s\)에 대해 \(\mR\)을 부분체 \(\Omega(z_1,\ldots,z_s)\)의 대수적 확대체로 나타낼 수 있으면 반드시 \(s=k\)이다.\footnote{임의의 표수에서 성립하는 이 알려진 사실의 증명은 Steinitz, § 22, 9를 보라. \(u\)들을 첨가해도 대수적 랭크는 커질 수 없다. 원래 원소들과 \(\Omega(u)\)의 계수로 이루어진 유리적 결합만을 다루기 때문이다. 또 작아질 수도 없다. \(\mR\)의 원소들 사이의 \(\Omega(u)\)-계수 관계는 유한개의 \(\Omega\)-계수 관계로 분해되기 때문이다.} 마찬가지로 \(u\)들이 \(\mR\)에 들어 있지 않은 부정원들이면, \(\Omega(u)\)에 관한 \(\mR(u)\)의 대수적 랭크도 \(k\)이다.

이제 한 부정원의 소함수에 통상적으로 쓰는 방법\footnote{Steinitz, § 6.}과 완전히 나란하게 다음 정리에 따라 영점체를 구성한다.

\textbf{정리 III.} \emph{차원 \(n-i\)인 소아이디얼 \(\frakp\)의 잉여류체 \((\mR)\)에는 그와 동형인 \(\frakp\)의 영점체 \(R\)을 대응시킬 수 있다. 이 체의 대수적 랭크는 \(n-i\), 곧 \(n-i\)개의 매개변수에 의존하며, 특히 \(P(u,x_{i+1},\ldots,x_n)\)의 유한확대체로 볼 수 있다. \((\mR)\)과 \(R\) 사이의 동형은 \((P)\)와 \(P\), \((P(u))\)와 \(P(u)\) 사이의 동형을 포함하고, \(x_{i+1},\ldots,x_n\)을 매개변수로 삼으면 \((P(u,x_{i+1},\ldots,x_n))\)와 \(P(u,x_{i+1},\ldots,x_n)\) 사이의 동형도 포함한다. 역으로 대수적 랭크가 \(n-i\), 곧 \(n-i\)개의 매개변수에 의존하는 모든 영점체는 잉여류체 \((\mR)\)과 동형이다.}

정리 III의 따름정리로 알려진 다음 정리를 얻는다. \emph{한 아이디얼 \(\fraka\)가 소아이디얼 \(\frakp\)의 모든 영점에서 사라지면, \(\fraka\)는 \(\frakp\)로 나누어진다.}

증명을 위해 먼저 \((P(u))\)에 관한 잉여류체 \((\mR)\)의 대수적 랭크가 \(n-i\)임을 주목한다. \((x_{i+1}),\ldots,(x_n)\)의 류들은 \((P(u))\)에 관하여 대수적으로 독립이다. 차원이 \(n-i\)인 \(\frakp\)는 \(x_1,\ldots,x_i\)와 무관한 다항식을 포함하지 않기 때문이다. 그러나 다른 모든 류는 이 류들에 종속이다. 초등인자형 \(P^{(i)}\)이 \(\frakp\)에 속하므로 § 1, 2에 따라 \(\lambda=1,2,\ldots,i\) 각각에 대하여 \(x_\lambda,x_{i+1},\ldots,x_n\)에 의존하는 다항식도 \(\frakp\)에 속하기 때문이다. 따라서 \((\mR)\)은 \((P(u,x_{i+1},\ldots,x_n))\)의 유한확대체가 된다. § 1, 5에서 \(\frakg_{i-1}\)은 단위아이디얼이고 정리 I에 따라 \(\frakg_i=\frakp\)임을 고려하면, 이 부분체에 관한 차수는 노름의 차수와 같다.

동형인 영점체로 넘어가기 위해서는 \(P(u,x_{i+1},\ldots,x_n)\)과 \((P(u,x_{i+1},\ldots,x_n))\) 사이에 다음 동형이 있음을 더 주목해야 한다. \(P(u,x_{i+1},\ldots,x_n)\)의 각 원소에 그 원소가 나타내는 잉여류를 대응시키면 된다. 우선 \(x_{i+1},\ldots,x_n\), 또는 \((x_{i+1}),\ldots,(x_n)\)의 다항식들로 이루어진 정역들이 이 대응으로 동형이다. \(\frakp\)에는 \(x_1,\ldots,x_i\)와 무관한 다항식이 없으므로 \(P(u,x_{i+1},\ldots,x_n)\)의 서로 다른 다항식들에는 서로 다른 류들이 대응하기 때문이다. 동형인 정역들에서 분수몫을 만들면 동형인 체들이 생긴다. \(P(u,x_{i+1},\ldots,x_n)\)과 \((P(u,x_{i+1},\ldots,x_n))\) 사이의 이 동형은 \(P\)와 \((P)\), \(P(u)\)와 \((P(u))\) 사이의 동형을 포함한다.

이 동형을 \((\mR)\)과 한 영점체 \(R\) 사이의 동형으로 확대하기 위하여, 류 \((x_1),\ldots,(x_i)\)에 어떤 새 원소 \(\bar{x}_1,\ldots,\bar{x}_i\)를 동형으로 대응시키자. 곧 각 다항식 \((a(x))\)에는 다항식 \(a(\bar{x}_1,\ldots,\bar{x}_i,x_{i+1},\ldots,x_n)\)이 대응하고, 합에는 합이, 곱에는 곱이 대응하게 한다. 이제 분수몫을 만들면 --- 분모를 각각 부분체 \((P(u,x_{i+1},\ldots,x_n))\)와 \(P(u,x_{i+1},\ldots,x_n)\)의 원소들로 제한해도 된다 --- 잉여류체 \((\mR)\)에 체 \(R\)이 대응한다. \(R\)은 \(P(u,x_{i+1},\ldots,x_n)\)의 노름 차수의 유한확대체이며, 영점체라고 부를 수 있다. 실제로 이 대응은 \(f(x)\equiv0(\frakp)\)이면 \(f(\bar{x}_1,\ldots,\bar{x}_i,x_{i+1},\ldots,x_n)\)이 사라진다고 말한다. 이 전체 논의에서 \((P(u,x_{i+1},\ldots,x_n))\) 대신 \((\mR)\)의 대수적으로 독립인 \(n-i\)개 원소를 첨가하여 생기는 다른 임의의 부분체를 써도 된다.

역으로 대수적 랭크가 \(n-i\)인 모든 영점체 \(R\)이 잉여류체와 동형임을 쉽게 보일 수 있다. \(R\)은 \(\bar{x}_1,\ldots,\bar{x}_n\)에서 \(P(u)\)의 계수를 갖는 유리식으로 얻어지므로, 이 원소들 가운데 \(n-i\)개는 대수적으로 독립, 곧 \(P(u)\)에 관하여 초월적이어야 한다. 특히 \(\bar{x}_{i+1},\ldots,\bar{x}_n\)이 그러해야 한다. \(P^{(i)}\equiv0(\frakp)\)에서 앞에서처럼 \(R\)이 \(P(u,\bar{x}_{i+1},\ldots,\bar{x}_n)\)의 대수적 확대체임이 따르기 때문이다. 가정상 \(\frakp\)의 모든 다항식 \(f(x)\)에 대하여 \(f(\bar{x})\)가 사라지므로, \((\mR)\) 안의 \((x)\)들의 다항식영역에 속하는 각 류에는 \(\bar{x}\)들의 다항식인 \(R\)의 원소가 하나씩만 대응하고, 합에는 합이, 곱에는 곱이 대응한다.

서로 다른 류들에는 \(R\)의 서로 다른 원소들이 대응한다. 그렇지 않으면 영과 다른 어떤 류, 따라서 적당히 곱한 뒤에는 \((P(u,x_{i+1},\ldots,x_n))\)의 다항식이 나타내는 어떤 류에 \(R\)의 영원소가 대응하게 되어, 가정에 어긋나게 \(\bar{x}_{i+1},\ldots,\bar{x}_n\) 사이에 대수적 종속관계가 생긴다. 이 대응은 \(R\)에 들어 있는 \(\bar{x}\)들의 모든 다항식을 소진한다. 이제 분수몫을 만들면 --- 다시 분모를 각각 \((P(u,x_{i+1},\ldots,x_n))\)과 \(P(u,\bar{x}_{i+1},\ldots,\bar{x}_n)\)의 원소들로 제한해도 된다 --- 이 동형인 정역들에서 동형인 체들이 생긴다.

이제 아이디얼 \(\fraka\)가 \(\frakp\)의 모든 영점, 특히 \(n-i\)개의 매개변수에 의존하는 모든 영점에서 사라진다고 하자. \(a(x)\)가 \(\fraka\)의 임의의 다항식이면 \(a(\bar{x})\)가 사라진다. \(R\)과 \((\mR)\) 사이의 동형에 따라 \((a(x))\)은 영류이고, 따라서 \(a(x)\), 나아가 \(\fraka\)가 \(\frakp\)로 나누어진다. 이로써 정리 III과 따름정리를 증명하였다.

정리 III에서 § 1, 9에 준 분해로 넘어가려면, 우선 소아이디얼의 경우에도 지수에 관하여 더 정확한 명제를 얻기 위해 초등인자형을 더 조사해야 한다.

\textbf{보조정리 V.} \emph{표수가 \(p\)일 때 소아이디얼 또는 으뜸아이디얼의 초등인자형 \(E^{(i)}\), 더 일반적으로 최고차원과 단순한 차원이 일치하는 아이디얼 \(\fraka\)의 초등인자형이 \(x_i\)에 관하여 지수 \(f\), 곧 \(x_i^{p^f}\)의 다항식이면, \(x_{i+1},\ldots,x_n\)과 \(E^{(i)}\)에 나타나는 부정원 \(u_{\mu r}\), 또는 \(t_{\mu r}\)에 관해서도 지수 \(f\)이다. 특히 계수영역이 완전체이면 소아이디얼의 초등인자형 \(P^{(i)}\)은 언제나 \(x_i\)에 관하여 지수 0, 곧 제1종 소함수이다.}

여기서 \(P\)가 표수 \(p\)의 완전체이면 \(P(u,x_{i+1},\ldots,x_n)\)은 비완전체가 되므로, § 1, 14만으로는 \(P^{(i)}\)이 제1종이라고 곧바로 결론 내릴 수 없다는 점을 주의해야 한다. 이제 \(E^{(i)}\)이 \(x_i\)에 관하여 지수 \(f\)이지만 다른 부정원 \(x_{i+\lambda}\)를 지수 \(f'\)로 포함한다고 하자. \(u_{\mu r}\), 또는 \(t_{\mu r}\)들을 서로 바꾸면 § 1, 2에 따라 아이디얼로 나누어지는 영과 다른 다항식 \(G^{(i)}\)도 존재한다. 이는 \(x_i^{p^{f'}}\)의 다항식이고, 초등인자형의 정의에 따라 \(E^{(i)}\)으로 나누어진다. 그러나 \(G^{(i)}\)과 \(E^{(i)}\)의 차수가 같으므로, \(G^{(i)}\)은 \(P(u)\)의 한 원소를 인자로 한 것만 제외하면 \(E^{(i)}\)과 일치해야 하고, 따라서 반드시 \(f'=f\)이다. 그러므로 \(t_{\mu r}\)들에 관하여 다항식이고 원시적이라고 가정할 수 있는 \(E^{(i)}\)은 다음 꼴이다.
\[
  E^{(i)}=\sum c_\lambda(t)\,x_i^{\lambda_i p^f}\cdots x_n^{\lambda_n p^f}
       =\sum c_i(t)\,g_i(y^{p^f},t^{p^f})
       =\sum T_\mu(t)\,\bar h_\mu(y^{p^f}),
\]
여기서 \(T_\mu\)는 \(t\)들의 거듭제곱곱이고, \(\bar h\)들은 \(\bar{\fraka}\)로 나누어진다. 따라서 \(T_\mu\)에서 각 \(t\)의 지수를 그 안에 포함된 \(p^f\)의 최대배수로 바꾸어도 \(\fraka\)에 속하는 다항식을 얻는다. 위 표현에서 보듯 이는 \(c_\lambda(t)\)에서 각 \(t\)의 지수를 그 안에 포함된 \(p^f\)의 최대배수로 바꾸는 것과 같다. 이 치환으로 \(E^{(i)}\)은 \(t,x_i,\ldots,x_n\)의 한 다항식으로 가는데, 정의에 따라 이 다항식은 \(E^{(i)}\)으로 나누어지며, 가정한 원시성 때문에 \(t\)들에 관해서도 그러하다. 어느 인수에 관해서도 차수가 늘지 않았으므로 이 다항식은 \(E^{(i)}\)과 동일해야 한다. 특히 \(P\)가 완전체이면 \(E^{(i)}\)은 \(x^{p^f},t^{p^f}\)의 다항식으로서 \(p^f\)제곱이다. 따라서 소아이디얼을 다룰 때에는 반드시 \(f=0\)이고, 초등인자형 \(P^{(i)}\)은 제1종 소함수이다.

이제 분해는 다음 정리로 주어진다.

\textbf{정리 IV.} \emph{소아이디얼의 초등인자형 \(P^{(i)}\)은 영점체에서 유도된 갈루아체 안에서 다음과 같이 명시적으로 분해된다.}
\[
 P^{(i)}=\prod (x_i-t_{1i}\bar y_{1\nu}-\cdots-t_{ni}\bar y_{n\nu})^\delta
 =\prod (t_{1i}(y_1-\bar y_{1\nu})+\cdots+t_{ni}(y_n-\bar y_{n\nu}))^\delta,
\]
\emph{여기서 \(\bar y_{1\nu},\ldots,\bar y_{n\nu}\)는 각각 서로 대응하고 \(t_{1i},\ldots,t_{ni}\)와 무관한 원래 부정원 \(y\)들의 영점계이며, 서로 다른 \(\nu\)에는 서로 다른 일차인자가 대응한다. 지수 \(\delta\)는 \(P\)가 완전체이면 1이고, 비완전체이면 \(p^f\) (\(f\geq0\))이다. 이로써 새 부정원 \(v\)들을 첨가하고 \(u_{\mu\nu}\) 대신 \(u\)와 \(v\)의 쌍선형 결합을 넣어 쓴 초등인자형의 분해도 함께 정해진다.}
\[
 P^{(i)}(u,v)=\prod (z-v_1\bar x_{1\nu}-\cdots-v_i\bar x_{i\nu})^\delta,
\]
\emph{여기서 \(\delta\)는 위와 같은 뜻이다. 마찬가지로 소아이디얼의 노름, 또는 \(\frakp\)를 대응 소아이디얼로 갖는 으뜸아이디얼의 노름이 \(P^{(i)}\)의 거듭제곱으로 분해되는 것도 함께 주어진다.}

정리 IV의 따름정리는 다음과 같다. \emph{두 소아이디얼의 초등인자형 \(P^{(i)}\)이 일치하면 두 소아이디얼은 동일하다.}

정리 IV의 증명은 초등인자형 \(P^{(i)}\)이 \((P(u,x_i,x_{i+1},\ldots,x_n))\)을 \((P(u,x_{i+1},\ldots,x_n))\) 위에서 본 확대체의 기본방정식과 동일함을 보이는 것으로 귀착된다. 먼저 부정원 \(u_{\mu r}\), 또는 \(t_{\mu r}\)에 대한 의존성을 알아보기 위하여 \(\bar{\frakp}\)의 잉여류체 \((\bar{\mR})\)로 넘어가자. 대수적 랭크의 정의에서 말한 바에 따라 이 체는 \((\bar P)\)에 관하여 대수적 랭크 \(n-i\)이어야 한다. 여기서 \((\bar P)\)는 \(P\)의 원소들이 나타내는 부분체이다. 실제로 \((\bar{\mR})\)은 \((\mR)\)과 동형인 한 부분체에 부정원 \(u\), 또는 \(t\)를 첨가하여 생긴다. 이 부분체는 \(\bar{\frakS}\)의 다항식들로 나타낼 수 있는 모든 류와 오직 그 류들에서 분수몫을 만들어 얻으며, \((\bar P)\)와 \((P)\)가 서로 대응한다. 따라서 \((P(u))\)에 관한 \((\bar{\mR})\)의 대수적 랭크는 \((\bar P)\)에 관한 이 부분체의 대수적 랭크와 같고, 동형에 따라 다시 \((P)\)에 관한 \((\mR)\)의 대수적 랭크와 같다. \((\bar{\mR})\)은 류 \((y_1),\ldots,(y_n)\)에서 \((\bar P)\)의 계수를 갖는 유리식으로 얻어지므로 이 류들 가운데 \(n-i\)개는 대수적으로 독립이어야 하며, \((\bar{\mR})\)은 이 \(n-i\)개 류를 \((\bar P)\)에 첨가하여 생기는 부분체의 유한 대수적 확대체이다.\footnote{이 관찰은 정리 III과 그 따름정리를 \(\bar{\frakp}\)의 영점체 \(\bar R\)에 대하여 직접 진술하고 증명할 수도 있었음을 보여 준다. 그러나 변환된 아이디얼로 넘어가야 임의의 아이디얼에서 같은 차원을 가진 최대 으뜸성분들의 영점에 같은 매개변수들이 나타난다. 그래야 비로소 임의의 아이디얼의 노름과 초등인자형과의 연결이 생긴다.}

\(x_{i+1},\ldots,x_n\)에 나타나는 부정원 \(t_{\sigma\tau}\)들만 첨가하면, 류 \((y_1),\ldots,(y_n)\)과 \((x_{i+1}),\ldots,(x_n)\)을 포함하고 \((P(t_{\sigma\tau},x_{i+1},\ldots,x_n))\)의 유한 대수적 확대체인 또 하나의 잉여류체를 얻는다. 여기의 체들은 \((\bar{\mR})\)의 부분체들과 동형이지만 동일하지는 않다. 부정원들이 다르면 잉여류들이 같은 원소들로 나타나더라도 같은 원소들을 포함하지 않아 동일하지 않기 때문이다. 따라서 \((P(t_{\sigma\tau},x_{i+1},\ldots,x_n))\)에 대한 류 \((y)\)들의 의존성에는 부정원 \(t_{\sigma\tau}\)들만 들어간다.

이제 \(t_{1i},\ldots,t_{ni}\)를 첨가하여 류 \((x_i)=(t_{1i}y_1+\cdots+t_{ni}y_n)\)을 만들고,
\(
 (x_{i\nu})=(t_{1i}y_{1\nu}+\cdots+t_{ni}y_{n\nu})
\)
를 \((P(t_{\sigma\tau},t_{1i},\ldots,t_{ni},x_{i+1},\ldots,x_n))\)에 관한 갈루아체 안의 서로 다른 \(r\)개 공액값이라 하자. 확대가 제1종인지 아닌지에 따라 인자 \(t-(x_{i\nu})\)들의 곱, 또는 그 \(p^f\)제곱들의 곱은 \((P(t,x_{i+1},\ldots,x_n))\)에 관한 소함수 \((G)\)이다. 실제로 차수가 \(r\cdot p^f\)인 원소들이 존재하므로(§ 1, 15), 각 원소가 \(t_{1i},\ldots,t_{ni}\)의 특수화로 생기는 이상 \((x_i)\)도 반드시 그런 원소이다. 그리고 \((x_i)\)는 \((G)\)의 영점이다. \(\frakp\)의 동형인 영점체와 거기서 유도된 갈루아체로 넘어가면, 여기서 \(G\)는 일차인자 \(t-t_{1i}\bar y_{1\nu}-\cdots-t_{ni}\bar y_{n\nu}\)들, 또는 그 \(p^f\)제곱들로 분해된다. 한편 \((G)\)는 \(t=(x_i)\)에서 사라지므로 \(G(x_i)\)은 \(\frakp\), 따라서 \(P^{(i)}\)으로 나누어진다. \(G(x_i)\)은 \(P(u,x_{i+1},\ldots,x_n)\)에 관한 소함수이고 \(x_{i+1},\ldots,x_n\)에 관하여 원시적이므로 \(P^{(i)}\)과 동일하다. 이로써 정리 IV의 첫 분해를 증명하였다.

둘째 분해로 넘어가기 위하여, 부정원을 첨가해도 제1종 확대인지 아닌지라는 성질은 유지됨을 주목하자. \(v_1,\ldots,v_i\)와 모든 \(t_{\mu\nu}\)를 첨가한 뒤 류 \((z)=(v_1x_1+\cdots+v_ix_i)\)와 공액류 \((z_\nu)=(v_1x_{1\nu}+\cdots+v_ix_{i\nu})\)를 만들면, 모든 \(t-(z_\nu)\)의 곱, 또는 그 \(p^f\)제곱들의 곱은 다시 \(t=(z)\)에서 사라지는 소함수 \((H)\)이다. 따라서 \(H\)는 위와 같이 일차인자들의 곱으로 나타낼 수 있고, 변환 (1)에 \(z=v_1x_1+\cdots+v_ix_i\)를 합성하여 \(\frakp\)에서 얻은 소아이디얼의 초등인자형과 동일하다. 상호 특수화 때문에 이 초등인자형은 \(P^{(i)}\)에서 \(u_{\mu\nu}\)들을 \(u\)와 \(v\)의 어떤 쌍선형 결합들로 바꾸어 얻어진다.\footnote{H.-N., § 7을 보라.} 이 분해들로써 노름이 \(P^{(i)}\)의 한 거듭제곱으로 분해되는 것도 주어진다.\footnote{제1종 확대, 특히 완전체에서는 첫째 거듭제곱이고, 비완전체에서는 \(p^g\)제곱이다. 이는 § 6의 정리 XIII과 XIV에서 보일 것이다.} 마찬가지로 \(\frakp\)를 대응 소아이디얼로 갖는 으뜸아이디얼의 노름과 초등인자형의 분해도 주어진다.

끝으로 두 소아이디얼의 초등인자형 \(P^{(i)}\)이 일치하면 위 분해에 따라 그 영점들도 일치한다. 따라서 각 아이디얼은 다른 아이디얼로 나누어지고, 두 아이디얼은 동일하다.

\subsection*{§ 4. 차원개념의 산술적 정식화}

변환된 아이디얼의 도입과 무관한 차원개념의 첫 정식화는 다음 정리가 준다.

\textbf{정리 V.} \emph{한 소아이디얼의 차원은 그 잉여류체의 대수적 랭크로 주어진다. 한 아이디얼의 최고차원은 그 대응 소아이디얼들의 차원수 가운데 가장 큰 것과 같다.}

정리 V의 첫 부분은 정리 IV의 첫째 주에서 이미 증명하였다. 거기서는 \(\bar{\frakp}\)의 잉여류체 \((\bar{\mR})\)의 대수적 랭크가 \(\frakp\)의 잉여류체 \((\mR)\)의 대수적 랭크와 같고, 따라서 § 1, 8과 일치하게 \(\frakp\)의 차원과 같음을 보였다. 둘째 부분을 증명하기 위해 \(\frakm=[\frakq_1,\ldots,\frakq_a]\)가 극대 으뜸성분들에 의한 최단표현이라고 하자. 우선 각 \(\frakq\)는 \(\frakm\)의 인자이므로 \(\frakm\)보다 높은 차원을 가질 수 없다. 한편 모든 \(\frakq\)의 곱은 \(\frakm\)으로 나누어지고, 따라서 개별 \(\frakq\)의 모든 초등인자형의 곱도 \(\frakm\)으로 나누어진다. 그러므로 \((n-i)\)가 \(\frakq\)들의 차원수 가운데 가장 크다면 \(\frakm\)은 영과 다른 다항식 \(G^{(i)}\)을 포함하고, 따라서 그 차원은 \(n-i\)보다 높을 수 없다. 정리 I에 따라 \(\frakq\)들의 차원수는 그 대응 소아이디얼들의 차원수와 일치한다. 변환된 아이디얼을 도입하지 않고 소아이디얼의 차원을 잉여류체의 대수적 랭크로 정의하면, 정리 V의 둘째 부분은 임의의 아이디얼의 최고차원을 정의한다.

한편 소아이디얼들의 사슬을 이용하여 역시 변환된 아이디얼의 도입과 무관하게 차원개념을 정식화하려면, 정리 II를 보충하여 다음을 보여야 한다.

\textbf{정리 VI.} \emph{단위아이디얼과 다른 모든 소아이디얼은 --- 필요하면 부정원들을 첨가한 뒤 --- 차원이 정확히 하나 작은 소아이디얼을 적어도 하나 인자로 갖는다.}

실제로 \(\frakp\)의 차원이 \((n-i)>0\)이고 \(v\)가 새 부정원이면
\[
 \frakc=(\frakp,x_i-v)
\]
는 원하는 성질을 가진 아이디얼이다. 이는 동형대응에서 따른다. 곧 \(\frakc=(\frakp^*,x_i-v)\)이며, 여기서 \(\frakp^*\)는 \(\frakp\)에서 \(x_i\)를 부정원 \(v\)로 바꾸고 \(v\)를 계수영역에 첨가하여 얻는다. 잉여류들도 마찬가지로 류 \((x_i)\)를 \((v)\)로 바꾸어 얻는다. 그런데 \(x_i\)를 \(P(u)\)에 첨가해도 \(\frakp\)는 자기 자신으로 간다. 실제로
\[
 h(x_i)f(x)\equiv0(\frakp),\qquad h(x_i)\ne0
\]
이면 반드시 \(f(x)\equiv0(\frakp)\)이다. \(\frakp\)가 \(x_i\)에만 의존하는 다항식을 포함한다면 \(x_n\)에만 의존하는 다항식도 포함하여, 가정에 어긋나게 차원이 기껏해야 0이 되기 때문이다. 따라서 대응 \(v\sim x_i\) 아래에서 \(\frakc\)에 따른 영류에는 반드시 \(\frakp\)에 따른 영류가 대응하고, 그 역도 물론 성립한다. 이로써 \(\frakp\)와 \(\frakc\)에 따른 잉여류계들 사이에 동형대응이 생긴다. 뒤의 잉여류계에는 영인자가 없으므로 \(\frakc\)는 소아이디얼이다.

이 동형대응 아래에서 \((x_i)\)와 \((P(u,x_{i+1},\ldots,x_n))\) 사이의 대수적 종속관계에는 \((P(u,v))\)에 관한 \((x_{i+1}),\ldots,(x_n)\) 사이의 한 관계가 대응한다. 반면 \(\lambda=1,\ldots,i-1\)에 대하여 \((x_\lambda)\)와 \((P(u,x_{i+1},\ldots,x_n))\) 사이의 종속관계들은 그대로 유지되어 대수적 랭크를 더 낮추지 않는다. 따라서 대수적 랭크는 정확히 하나 작아졌다. \(\frakp\)의 차원이 0이었다면 \(\frakc\)는 단위아이디얼이 되므로 위 결과는 여전히 성립한다. 따라서 원래 영역 \(\bar{\frakS}\)에서는
\[
 \bar{\fraka}=(\bar{\frakp},v+v_1y_1+\cdots+v_ny_n),
\]
가 원하는 성질을 가진 아이디얼이다. 여기서 \(v_\lambda=-t_{\lambda i}\)로 두었다. 특히 \(P\)가 무한개의 원소를 포함하면 \(v,v_1,\ldots,v_n\)을 언제나 \(P\)의 원소들로 특수화하여, \(\bar{\fraka}\)의 노름과 초등인자형, 따라서 차원이 특수화된 아이디얼 \(\frakb\)의 것과 같게 할 수 있다.\footnote{H.-N., § 7의 마지막 단락을 보라.} 다만 이때 초등인자형은 더 이상 소함수로 남지 않을 수 있다. 정리 V에 따라 \(\frakb\)는 원하는 차원의 대응 소아이디얼을 적어도 하나 갖는다. 다시 \(\bar{\frakS}\)로 돌아가면, 이 경우 정리 VI는 어떠한 부정원도 첨가하지 않고 성립한다.

정리 II와 정리 VI는 곧바로 원하는 정식화를 준다.

\textbf{정리 VII.} \emph{한 소아이디얼 \(\frakp\)의 차원이 \(n-i\)라는 것은 --- 필요하면 부정원들을 첨가한 뒤 --- 다음과 같은 \(1+(n-i)+1\)개의 소아이디얼로 이루어진 사슬이 적어도 하나 존재한다는 뜻이다.}
\[
 \frakp_0=\frakp;\quad \frakp_1,\ldots,\frakp_{n-i};\quad \frakp_{n-i+1},
\]
\emph{여기서 각 아이디얼은 바로 앞 아이디얼의 진인자이고, 항의 수가 더 많은 그러한 사슬은 존재하지 않는다. \(P\)가 무한개의 원소를 포함하면 이 정의는 부정원을 전혀 도입하지 않고 원래 영역에서도 성립한다.}\footnote{차원의 이 사슬정의는 대수적 수체의 경우 통상적인 소아이디얼들에는 차원 0, 단위아이디얼에는 차원 \(-1\), 영아이디얼에는 차원 1을 준다. 실제로 영아이디얼도 소아이디얼의 정의를 만족한다. 체에서는 사라지지 않는 두 원소의 곱이 언제나 영과 다르기 때문이다. 이 차원정의 아래에서도 정리 II는 대수적 수체에서 성립한다.}

우선 사슬의 마지막 항은 소아이디얼의 정의를 만족하는 단위아이디얼이어야 한다. 그렇지 않으면 \(\frako\)를 덧붙여 사슬을 늘일 수 있기 때문이다. \(\frako\) 자체에는 정리 VII에 따라 § 1, 8과 일치하게 차원 \(-1\)이 주어진다. 이제 \(\frakp\ne\frako\)이고 그 차원이 \(n-i\)라 하자. 정리 II에 따라 같은 차원의 소아이디얼 두 개가 나타날 수 없으므로 사슬의 항 수는 기껏해야 \(1+(n-i)+1\)이다. 그러나 정리 VI에 따라 --- 필요하면 새 부정원들을 첨가하여 --- 각 항마다 차원이 정확히 하나씩 작아지는 이 항 수의 사슬이 적어도 하나 존재한다. 정리 VII를 차원의 정의로 삼으면 --- 이 정의는 원래 영역에서도 그대로 말할 수 있고, \(P\)가 무한개의 원소를 가지면 정리 VI에 따라 부정원을 도입할 필요도 없다 --- 임의의 아이디얼의 최고차원은 다시 정리 V의 둘째 부분으로 정의된다.

\subsection*{§ 5. 기초아이디얼과 그 성분들을 일반 아이디얼 이론 안에 배치하기. 아이디얼의 영점}

여기서는 기초아이디얼과 그 성분들(§ 1, 7)을 일반 아이디얼 이론의 뜻에서 분해의 고립성분들(§ 1, 13)로 특징짓는다. 이 특징은 노름의 분해정리와 일반 아이디얼 이론의 분해정리가 서로 나란히 진행됨을 보여 주고, § 3에서 소아이디얼과 으뜸아이디얼에 대하여 준 영점이론을 임의의 아이디얼에 대하여 진술할 수 있게 한다. 기초아이디얼에 대해서는 다음이 성립한다.

\textbf{정리 VIII.} \emph{\(i\)단계 기초아이디얼 \(\frakg_i\)는 차원이 \(n-1,n-2,\ldots,n-i\)인 모든 대응 소아이디얼들에 의해 유일하게 결정되는 분해의 고립성분이다. 모든 대응 소아이디얼의 차원이 같으면 아이디얼 자체도 그 차원을 갖는다. 곧 노름에는 인자 \(R^{(i)}\) 하나만 나타난다. 일반적으로 노름에는 대응 소아이디얼들 사이에 서로 다른 차원수가 나타나는 수만큼 단위원과 다른 인자 \(R^{(i)},R^{(i+\sigma)},R^{(i+\tau)},\ldots\)가 있다.}

증명에 앞서 다음을 둔다.

\textbf{보조정리 VI.} \emph{한 아이디얼이 최소공배수 \(\frakm=[\fraka_1,\ldots,\fraka_\alpha]\)로 나타나 있으면, 그 기초아이디얼 \(\frakg_i\)는 \(\fraka\)들의 \(i\)단계 기초아이디얼 \(\frakh_i\)들의 최소공배수이다.}

실제로
\[
 b^{(i+1)}f\equiv0(\frakm),\qquad b^{(i+1)}\ne0
\]
에서
\[
 b^{(i+1)}f\equiv0(\fraka_\lambda),\qquad b^{(i+1)}\ne0
\]
이 따른다. 따라서 \(\frakg_i\)는 \([\frakh_{i1},\ldots,\frakh_{i\alpha}]\)로 나누어진다. 역으로
\[
 c_\lambda^{(i+1)}f\equiv0(\fraka_\lambda),\qquad c_\lambda^{(i+1)}\ne0
\]
에서
\[
 c_1^{(i+1)}\cdots c_\alpha^{(i+1)}f\equiv0(\frakm),\qquad
 c_1^{(i+1)}\cdots c_\alpha^{(i+1)}\ne0
\]
도 따르므로 반대쪽 나눗셈도 성립한다. 이로써 보조정리를 증명하였다.

이제 \(\frakm\)의 대응 소아이디얼들을 차원에 따라 다음과 같이 배열하자. 물론 어떤 차원들은 빠질 수도 있으며, 그때 \(j_\lambda=0\)이다.
\[
 \frakp_{1,1},\ldots,\frakp_{1,j_1};\quad
 \frakp_{2,1},\ldots,\frakp_{2,j_2};\quad\ldots;\quad
 \frakp_{n,1},\ldots,\frakp_{n,j_n}.
\]
일반적으로 \(\frakp_{\lambda,\kappa}\)는 차원 \(n-\lambda\)인 소아이디얼을 뜻하고, \(\frakq_{\lambda,\kappa}\)는 \(\frakm\)의 어떤 고정된 분해에 나타나는 그에 대응하는 으뜸아이디얼이라 하자. 정의에 따라 \(\lambda>i\)인 모든 \(\frakq_{\lambda,\kappa}\)의 \(i\)단계 기초아이디얼은 단위아이디얼이고, \(\lambda\leq i\)이면 정리 I에 따라 그 기초아이디얼은 \(\frakq_{\lambda,\kappa}\)와 같다. 따라서 보조정리 VI에서
\[
 \tag{5}
 \frakg_i=[\frakq_{1,1},\ldots,\frakq_{1,j_1};\ldots;\frakq_{i,1},\ldots,\frakq_{i,j_i}]
\]
를 얻는다. 이 표현은 \(\frakg_i\)가 분해의 고립성분임을 보여 준다. \(\frakg_i\)에 속하는 어떤 소아이디얼도 모두 더 낮은 차원을 가진 나머지 소아이디얼 가운데 하나에 포함될 수 없기 때문이다.

특별히 대응 소아이디얼들 가운데 차원 \(n-i\)인 것들만 나타나면, (5)에 따라 \(\frakg_0,\ldots,\frakg_{i-1}\)은 단위아이디얼이고 \(\frakg_i=\frakg_{i+1}=\cdots=\frakm\)이다. 따라서 \(R_\frakm=R^{(i)}\)이다. 서로 다른 차원의 소아이디얼들이 있어서, 이를테면 \(j_i,j_{i+\sigma},j_{i+\tau},\ldots\)가 영과 다르면 (5)에 따라
\[
 \frakg_0\ne\frakg_i\ne\frakg_{i+\sigma}\ne\frakg_{i+\tau}\ne\cdots,
\]
이지만
\[
 \frakg_0=\frakg_1=\cdots=\frakg_{i-1}=\frako,
 \qquad
 \frakg_i=\frakg_{i+1}=\cdots=\frakg_{i+\sigma-1},\ldots
\]
이다. 따라서 인자 \(R^{(i)},R^{(i+\sigma)},R^{(i+\tau)}\)와 오직 이 인자들만 단위원과 다르다.\footnote{정리 VIII에서 보조정리 II도 함께 고려하면, 기초아이디얼은 변환된 아이디얼들의 최소공배수로서 다시 변환된 아이디얼임이 따른다. 정리 VIII은 이 사실을 쓰지 않는 H.-N.의 정리들에만 근거하므로, 이로써 내적인 이유까지 드러내는 새 증명을 얻는다. H.-N.에서는 이 정리가 영점이론에서만 쓰이며, 그 이론은 여기서 새로 전개하였다.}

정리 VIII을 이용하면 정리 I을 다음과 같이 강화할 수 있다.

\textbf{정리 IX.} \emph{한 아이디얼이 으뜸일 필요충분조건은 그 초등인자형, 따라서 그 노름이 으뜸함수인 것이다.}

으뜸아이디얼이 이 조건을 만족함은 정리 I에서 보였다. 역으로 \(\frakm\)의 초등인자형이 \(Q^{(i)}=(P^{(i)})^\rho\)이고 \(P^{(i)}\)이 소함수라고 하자. 정리 VIII에 따라 \(\frakm\)의 모든 소아이디얼은 같은 차원 \(n-i\)을 갖는다. \(\frakm=[\frakq_1,\ldots,\frakq_a]\)에서
\[
 Q^{(i)}\equiv0(\frakq_\lambda);\quad\hbox{따라서}\quad
 (P^{(i)})^\rho\equiv0(\frakp_\lambda)\quad\hbox{이고, 그러므로}\quad
 P^{(i)}\equiv0(\frakp_\lambda)
\]
를 얻는다. 이제 모든 \(\frakp_\lambda\)의 차원이 \(n-i\)이고 \(P^{(i)}\)은 소함수이므로, \(P^{(i)}\)은 각 \(\frakp_\lambda\)의 초등인자형과 같다. 따라서 정리 IV의 따름정리에 따라 모든 \(\frakp_\lambda\)가 일치한다. 극대 으뜸성분들을 다루고 있으므로 하나의 \(\frakp\)만 나타날 수 있고, \(\frakm\)은 으뜸이다. \(\rho=1\)일 때 소아이디얼이라는 특수한 경우도 배제되지 않는다. 정리 VIII과 IX에서 기초아이디얼의 성분들을 다음과 같이 배치할 수 있다.

\textbf{정리 X.} \emph{\(R^{(i)}\)의 극대 으뜸인자 \(Q_\lambda^{(i)}=(P_\lambda^{(i)})^{\rho_\lambda}\)에 대응하는 기초아이디얼 \(\frakg_i\)의 성분 \(\frakr_\lambda\)는, 차원이 \(n-1,\ldots,n-i+1\)인 대응 소아이디얼들과 초등인자형이 \(P_\lambda^{(i)}\)인 차원 \(n-i\)의 그 소아이디얼에 의해 유일하게 결정되는 분해의 고립성분이다. 대응 소아이디얼들과 노름의 극대 으뜸인자들은 일대일로 대응한다.}

§ 1, 7에 따라 \(\frakr_\lambda\)는 그 \(i\)단계 기초아이디얼과 일치하고, \(\frakg_{i-1}\)을 \((i-1)\)단계 기초아이디얼로 갖는다. 따라서 정리 VIII, 또는 (5)에 따라 다음 최단표현이 있다.
\[
 \frakr_\lambda=[\frakg_{i-1},\frakq_{\lambda,1},\ldots,\frakq_{\lambda,t_\lambda}],
\]
여기서 모든 \(\frakq\)의 차원은 \(n-i\)이고 서로 다른 소아이디얼들에 속한다. \(Q_\lambda^{(i)}\)의 정의와 \(\frakr_\lambda\)가 그 \(i\)단계 기초아이디얼과 일치한다는 사실에서
\[
 Q_\lambda^{(i)}\frakg_{i-1}\equiv0(\frakq_{\lambda,\kappa}),
 \quad(\kappa=1,\ldots,t_\lambda)
\]
를 얻고, 따라서
\[
\begin{aligned}
 (Q_\lambda^{(i)})^{\sigma_\lambda}&\equiv0(\frakq_{\lambda,\kappa});
 \quad\hbox{그러므로}\quad
 (P_\lambda^{(i)})^{\rho_\lambda\sigma_\lambda}\equiv0(\frakp_{\lambda,\kappa}),\\
 &\hbox{따라서}\quad P_\lambda^{(i)}\equiv0(\frakp_{\lambda,\kappa})
\end{aligned}
\]
이다. 정리 IX의 증명과 마찬가지로 여기서 \(\frakr_\lambda\)의 표현에는 차원 \(n-i\)인 \(\frakq_\lambda\) 하나와 그 대응 소아이디얼 \(\frakp_\lambda\) 하나만 나타날 수 있고, \(\frakp_\lambda\)의 초등인자형은 \(P_\lambda^{(i)}\)이다.

이제 \(\frakp_\lambda\)가 \(\frakm\)의 대응 소아이디얼임을, 곧 \(\frakq_\lambda\)가 \(\frakm\)을 극대 으뜸성분들로 나타내는 표현 가운데 적어도 하나에 실제로 나타남을 보여야 한다. 그러면 어떤 소아이디얼도 같은 차원이나 더 낮은 차원의 소아이디얼에 포함될 수 없으므로 \(\frakr_\lambda\)가 분해의 고립성분, 더 정확히는 \(\frakp_\lambda\)와 더 높은 차원의 모든 대응 소아이디얼들에 의해 정해지는 성분임도 드러난다. 이를 위해서는 \(\frakq_\lambda\)가 \(\frakg_i\)의 한 최단표현에 나타남만 보이면 된다. 정리 VIII에 따라 그러한 표현은 언제나 \(\frakm\)의 한 최단표현으로 보충할 수 있기 때문이다. 따라서
\[
 \frakg_i=[\frakr_1,\ldots,\frakr_\alpha]
          =[\frakg_{i-1},\frakq_1,\ldots,\frakq_\alpha]
\]
가 최단임만 보이면 된다. 이를테면 \(\frakq_\lambda\)가 그 보충부분에 포함되어
\(
\frakg_{i-1}\frakq_1\cdots\frakq_{\lambda-1}
\frakq_{\lambda+1}\cdots\frakq_\alpha
\)
가 \(\frakq_\lambda\)로 나누어진다고 하자. \(\frakg_{i-1}\not\equiv0(\frakq_\lambda)\)이므로 어떤 \(\mu\ne\lambda\)에 대하여 \(\frakq_\mu\)의 한 거듭제곱이 \(\frakq_\lambda\)로 나누어진다. 대응 소아이디얼 \(\frakp_\mu\)와 \(\frakp_\lambda\)의 차원이 같으므로 정리 II에 따라 이 소아이디얼들이 일치해야 한다. 그러나 정의상 그 초등인자형 \(P_\mu^{(i)}\)과 \(P_\lambda^{(i)}\)은 서로 다르므로 이는 불가능하다. 정리 VIII에 따라 차원 \(n-i\)인 모든 대응 소아이디얼이 \(\frakg_i\)에 대응하며, 위 표현은 각각의 그러한 소아이디얼에 노름의 극대 으뜸인자가 실제로 하나씩 대응함도 보여 준다. 이로써 일대일 대응을 증명하였다.

정리 X에서 정리 IV를 곧바로 다음과 같이 옮길 수 있다.

\textbf{정리 XI.} \emph{노름의 각 극대 으뜸인자 \(Q_\lambda^{(i)}=(P_\lambda^{(i)})^{\rho_\lambda}\)는 다음 곱표현을 갖는다.}
\[
 Q_\lambda^{(i)}
 =\prod (x_i-t_{1i}\bar y_{1\nu}-\cdots-t_{ni}\bar y_{n\nu})^{\delta_\lambda\rho_\lambda}
 =\prod (t_{1i}(y_1-\bar y_{1\nu})+\cdots+t_{ni}(y_n-\bar y_{n\nu}))^{\delta_\lambda\rho_\lambda},
\]
\emph{여기서 \(\bar y_{1\nu},\ldots,\bar y_{n\nu}\)는 각각 원래 부정원 \(y\)에서 대응 소아이디얼 \(\frakp_\lambda\)의 서로 대응하고 \(t_{1i},\ldots,t_{ni}\)와 무관한 영점계이고, \(\delta_\lambda\)는 \(1\), 또는 \(p^{f_\lambda}\)이다. 이로써 노름의 일차인수분해가 완전히 주어진다. \(Q_\lambda^{(i)}\)의 차수는 고립성분 \(\frakr_\lambda\)에 따른 잉여류 가운데 \(\frakg_{i-1}\)의 잉여류들로 나타나는 부분환이 분수몫으로 유도된 잉여류 부분체 \((P(u,x_{i+1},\ldots,x_n))\)에 관하여 갖는 차수를 준다. 으뜸아이디얼의 경우에는 모든 잉여류로 이루어진 환을 말한다. 새 부정원 \(v\)들을 첨가하면 다음 분해도 얻는다.}
\[
 Q_\lambda^{(i)}(u,v)
 =\prod (z-v_1\bar x_{1\nu}-\cdots-v_i\bar x_{i\nu})^{\delta_\lambda\rho_\lambda}.
\]

정리 X에 따라 \(P_\lambda^{(i)}\)이 \(\frakp_\lambda\)의 초등인자형임을 알았으므로, 실제로는 정리 IV의 분해를 대입했을 뿐이다. 차수에 관한 주도 § 1, 5와 7에서 말한 것을 달리 표현한 것일 뿐이다.\footnote{마찬가지로 H.-N.의 주 2)에 언급된 중복도의 정식화도 정리 X에서 곧바로 따른다. 실제로 \(x_{i+1},\ldots,x_n\)을 첨가하면 잉여류계 \(\frakg_{i-1}\mid\frakr_\lambda\)는 \(\frakt_\lambda\mid\frakr_\lambda\)와 동형이다. 여기서 \(\frakt_\lambda=[\frakg_{i-1},\frakr_1,\ldots,\frakr_{\lambda-1},\frakr_{\lambda+1},\ldots,\frakr_a]\)로 둔다. \(Q_\mu^{(i)}\)들이 서로 소임을 고려하면 선형형식의 가군으로 돌아가는 것(H.-N., 정리 V)에서 이 사실이 곧바로 드러난다. 또 \(\frakt_\lambda\mid\frakr_\lambda\)는 \(\frakt_\lambda\mid\frakq_\lambda\)와 동형이다. 이때 \(\frakt_\lambda\)는 \(\frakq_\lambda\)의 보충부분에 \(x_{i+1},\ldots,x_n\)을 첨가하여 생긴다.}

\subsection*{§ 6. 초등인자형과 노름에 의한 소아이디얼과 진으뜸아이디얼의 특징짓기}

정리 IX에서는 초등인자형과 노름으로 으뜸아이디얼을 특징지었지만, 소아이디얼을 으뜸아이디얼의 특수한 경우로 보는 데 그쳤다. 이제 소아이디얼과 진으뜸아이디얼을 분리해야 하며, 논의는 § 3과 나란히 진행된다. 필요조건이나 충분조건은 다음과 같이 쉽게 줄 수 있다.

\textbf{정리 XII.} \emph{소아이디얼의 필요조건은 초등인자형이 소함수인 것이고, 충분조건은 노름이 소함수인 것이다. 달리 말하면 소아이디얼의 초등인자형은 언제나 소함수이고, 진으뜸아이디얼의 노름은 언제나 진으뜸함수이다.}

소아이디얼의 초등인자형이 소함수임은 정리 I에서 이미 보였다. 이제 노름 \(R_\frakq\)가 소함수와 같다고 하자. 이 가정 아래에서 \(\frakq\)의 모든 진인자 \(\fraka\)가 더 낮은 최고차원을 가짐을 보이면, 정리 II에 따라 \(\frakq\)가 소아이디얼임이 드러난다. 실제로 § 1, 6에 따르면 노름 \(R_\fraka\)의 인자들 가운데 \(R_\frakq\)의 대응 인자와 처음으로 다른 것은 진인자이다. \(R_\frakq=P^{(i)}\)이므로 \(R_\fraka\)의 인자 \(R^{(i)}\)은 \(P^{(i)}\)의 진인자와 같아야 하고, 따라서 단위원이다. 그러므로 \(\fraka\)의 최고차원은 더 낮다.

둘째의 간단한 증명은 다음 정리에도 바탕이 되는 보조정리를 이용한다.

\textbf{보조정리 VII.} \emph{으뜸아이디얼의 초등인자형 \(E^{(i)}\), 더 일반적으로 차원 \(n-i\)의 대응 소아이디얼들만 갖는 한 아이디얼의 초등인자형에 따른 잉여류 가운데 다항식 \(H^{(i)}\)들로 나타나는 잉여류계는, 그 아이디얼에 따른 잉여류의 한 부분계와 동형이다. 초등인자형의 차수는 분수체 \((P(u,x_{i+1},\ldots,x_n))\)에 관한 이 잉여류계의 차수를 준다.}

실제로 같은 다항식 \(H^{(i)}\)으로 나타나는 류들을 서로 대응시키면 대응을 만들 수 있으며, 합에는 합이, 곱에는 곱이 대응한다. 이 대응은 일대일이기도 하다. 그 아이디얼에 따른 영류에 속하는 모든 다항식 \(H^{(i)}\)은 정의상 \(E^{(i)}\)으로 나누어지기 때문이다. 따라서 영류에는 \(E^{(i)}\)에 따른 영류가 대응하고, 그 역도 물론 성립한다.

이제 한 아이디얼의 노름이 소함수 \(P^{(i)}\)이고, 따라서 초등인자형과 같다고 하자. 그러면 초등인자형 \(P^{(i)}\)과 \(\frakq\)에 따른 잉여류계들의 \((P(u,x_{i+1},\ldots,x_n))\)에 관한 차수는 일치한다. 따라서 \(P^{(i)}\)에 따른 잉여류계와 동형인 부분계가 \(\frakq\)에 따른 잉여류들을 모두 소진한다. 이 계에는 영인자가 없으므로 \(\frakq\)는 소아이디얼이다.

일반적으로 정리 XII는 뒤의 예들이 보여 주듯 필요충분조건으로 강화할 수 없다. 그러나 계수영역 \(P\)가 완전체(§ 1, 14)이면 다음 정리가 성립한다.

\textbf{정리 XIII.} \emph{계수영역 \(P\)가 완전체이면 소아이디얼의 노름과 초등인자형은 소함수이고, 따라서 서로 같다. 진으뜸아이디얼의 노름과 초등인자형은 진으뜸함수이다. 그러므로 완전체에서는 초등인자형이 소함수라는 조건이 소아이디얼의 필요충분조건이다. 이 조건에서 초등인자형 대신 노름을 써도 된다.}

정리 XII를 고려하면, 완전체에서 초등인자형이 소함수인 순간 그 아이디얼이 소아이디얼이고 그 노름도 이 소함수와 같음을 보이면 정리 XIII가 증명된다. 보조정리 V, § 3에서 완전체의 경우 초등인자형이 소함수이기만 하면 제1종 소함수임을 보였다. 따라서 \((\mR)\)을 \((P(u,x_{i+1},\ldots,x_n))\)에 \((x_i)\)를 첨가하여 얻은 \(P^{(i)}\)에 따른 잉여류체 \((P(u,x_i,x_{i+1},\ldots,x_n))\)라고 하자. 그러면 \((x_i)\)는 제1종 소함수 \((P^{(i)})\)의 영점이다. 따라서 이를테면 라그랑주 공식이 보여 주듯 \((y_1),\ldots,(y_n)\)은 \((\mR)\)에 들어 있다.

보조정리 VII에 따라 \((\mR)\)과 동형인, 고찰하는 아이디얼 \(\frakq\)에 따른 잉여류의 부분계 --- 여기서는 \((P(u,x_{i+1},\ldots,x_n))\)의 원소들이 분모에 나타날 뿐이다 --- 는 잉여류 \((y_1),\ldots,(y_n)\)을 포함하고, 따라서 \(\frakq\)에 따른 잉여류계를 모두 소진한다. 이 부분계는 \((\mR)\)과 동형이어서 영인자를 가질 수 없으므로 \(\frakq\)는 소아이디얼이다. \((P(u,x_{i+1},\ldots,x_n))\)에 관한 이 \(\frakq\)의 잉여류체 \((\mR)\)의 차수는 \(\frakq\)의 노름 차수와 같다. 한편 \((\mR)\)과의 동형에 따라 초등인자형 \(P^{(i)}\)의 차수와도 같으므로, 노름과 초등인자형은 일치해야 한다.

비완전체의 경우에도 정리 XII를 다음과 같이 강화할 수 있다.

\textbf{정리 XIV.} \emph{계수영역 \(P\)가 표수 \(p\)의 비완전체이면, 한 소아이디얼의 노름은 그 초등인자형의 \(p^g\)제곱이다. 여기서 \(0\leq g\leq(i-1)f\)이고, \(f\)는 초등인자형의 지수, \(n-i\)는 소아이디얼의 차원이다.}

정리 XIV는 \(\frakp\)의 잉여류체 \((\mR)\)이 그와 동형인 부분체
\(
(\mR')=(P(u,x_i,x_{i+1},\ldots,x_n))
\)
에 관하여 \(p^g\)차 (\(g\geq0\))라는 주장으로 귀착된다. 이는 § 1, 15에 제시한 Steinitz의 정리들, 더 정확히는 다음 보조정리에서 곧바로 따른다.

\textbf{보조정리 VIII.} \emph{\(\mathfrak L\)이 비완전체 \(\Omega\)의 축약차수 \(r\), 지수 \(f\)인 유한확대이고, \(\mathfrak L_0\)가 \(\mathfrak L\) 안의 제1종 부분체이며, \(\mathfrak M\)이 \(\mathfrak L_0\)와 \(\mathfrak L\) 사이의 중간체이면, \(\mathfrak L\)의 각 원소가 \(\mathfrak M\)에 관하여 갖는 차수는 \(p\)의 거듭제곱이다. 여기서 \(p\)는 \(\Omega\)의 표수이다.}

실제로 \(\mathfrak L\)의 각 원소 \(z\)의 \(p^{f'}\)제곱 (\(f'\leq f\))은 \(\mathfrak L_0\)에 속한다. 따라서 \(z\)는 \(\mathfrak L_0\)의 계수를 갖는 소함수 \(G(t)=(t-z)^{p^{f'}}\)의 영점이다. \(H(t)=(t-z)^\delta\)를 \(z\)를 영점으로 갖는 \(\mathfrak M\) 위의 소함수라 하자. 그러면 \(H(t)\)는 \(\mathfrak L_0(z)\)의 계수도 가지므로, 그 계수들은 \(\mathfrak M\)과 \(\mathfrak L_0(z)\)의 교차체, 곧 \(\mathfrak L_0\)와 \(\mathfrak L_0(z)\) 사이의 한 중간체에 속한다. 따라서 \(\mathfrak M\)에 관한 \(z\)의 차수는 이 중간체에 관한 차수와 같고, \(p^{f'}\)의 약수이므로 \(p\)의 거듭제곱이다.

정리 XIV를 증명하려면 다음만 주목하면 된다. \((\mR)\)의 축약차수가 \(r\), 지수가 \(f\)이면 체
\(
(\mR')=(P(u,x_i,x_{i+1},\ldots,x_n))
\)
는 \((P(u,x_{i+1},\ldots,x_n))\)에 관하여 \(r\cdot p^f\)차이고, 따라서 \((\mR)\) 안의 제1종 부분체 \((\mR_0)\)는 \((\mR')\)의 부분체이다. 실제로 \((\mR)\) 안에는 차수 \(r\cdot p^f\)인 원소들이 존재하므로(§ 1, 15), \((\mR)\)의 모든 원소가 \((x_i)\)에서 \(t_{\mu\nu}\)들을 특수화하여 생기는 이상 \((x_i)\)도 반드시 이 차수를 갖는다. 따라서 \((\mR)\)의 지수는 초등인자형의 지수와 일치하고, \((x_i)^{p^f}\)는 차수 \(r\)이며 \((\mR_0)\)의 원소이므로 \((\mR_0)\)의 원시원소이다. 이제 \((\mR)\)은 유한개의 원소, 이를테면 \((x_1),\ldots,(x_{i-1})\)을 \((\mR')\)에 첨가하여 생기므로 보조정리 VIII을 유한번 적용하면 원하는 증명과 동시에 \(g\)의 상계를 얻는다.\footnote{이와 나란히 진행되는 A. Ostrowski, \emph{Zur arithmetischen Theorie der algebraischen Größen}, Gött. Nachr. 1919, 279--298쪽, 특히 대응 정리가 증명 없이 진술된 288쪽을 보라.}

끝으로 정리 XII를 일반적으로 정리 XIV보다 더 강화할 수 없음을 보여 주는 예들을 덧붙인다. 노름이 초등인자형의 제곱인 소아이디얼 (\(g>0\))과, 초등인자형이 소함수인 진으뜸아이디얼들을 제시할 것이다. 뒤의 예는 여기에는 정리 XIV의 유사물이 없고, 노름이 초등인자형의 임의의 거듭제곱이 될 수 있음도 함께 보여 준다.

계수영역 \(P\)를 법 \(2\)의 잉여류체에 두 부정원 \(\lambda,\mu\)를 첨가하여 얻자. 이는 표수 2의 비완전체이다. 아이디얼
\[
 \bar{\frakp}=(y_1^2+\lambda,\,y_2^2+\mu)
\]
를 고찰하자. 그 잉여류계 \((\bar{\mR})\)이 체 \(P(\sqrt\lambda,\sqrt\mu)\)와 동형이므로 이는 소아이디얼이다. \((\bar{\mR})\)은 \((P)\)에 관하여 4차, 축약차수 \(r=1\), 지수 \(f=2\)이다. 따라서 초등인자형 \(P^{(2)}\)은 2차이고 노름은 4차여서 \(P^{(2)}\)의 제곱이다. § 1, 5의 모듈표현 (3)으로도 이를 쉽게 검산할 수 있으며, 다음을 얻는다.
\[
 P^{(2)}=(u_{11}u_{22}+u_{12}u_{21})^2x_2^2+(u_{11}^2\mu+u_{21}^2\lambda),
\]
또는 적당히 곱하면
\[
 x_2^2+(t_{12}^2\lambda+t_{22}^2\mu).
\]

이제 \(P\)를 법 \(2\)의 잉여류체에 부정원 \(\lambda\) 하나를 첨가하여 얻고, 아이디얼
\[
 \bar{\frakq}=(y_1^2+\lambda,\,y_2^2+\lambda)
              =(y_1^2+\lambda,(y_1+y_2)^2)
\]
를 바탕으로 삼자. 위에 준 초등인자형에서 \(\mu=\lambda\)로 특수화하면
\[
 P^{(2)}=x_2^2+\lambda(t_{12}^2+t_{22}^2)
\]
를 얻는다. \(t_{12}=(0)\), \(t_{22}=(1)\)로 두면 소함수 \(x_2^2+\lambda\)가 되므로 이것은 소함수이다. 그러나 \(\bar{\frakq}\)는 \((y_1+y_2)^2\)을 포함하지만 \((y_1+y_2)\)는 포함하지 않으므로 진으뜸이다. \(\bar{\frakq}\)에 따라 \(P\)에 관하여 선형독립인 잉여류가 네 개이므로, 노름은 \(P^{(2)}\)의 제곱, 곧 \(p\)제곱과 같아야 한다.

정리 XIV의 이 마지막 유사물이 언제나 성립하지는 않음을 다음 예가 보여 준다. \(P\)를 법 \(3\)의 잉여류체에 부정원 \(\lambda\)를 첨가하여 얻고,
\[
 \bar{\frakq}=(y_1^3+\lambda,(y_1-y_2)^2)
\]
를 바탕으로 삼자. 그러면 위와 같이 진으뜸아이디얼이다. \(y_2^3+\lambda\)도 \(\bar{\frakq}\)로 나누어짐을 고려하면 초등인자형은
\[
 P^{(2)}=x_2^3+\lambda(t_{12}+t_{22})^3
\]
이므로 소함수이다. 그러나 \(\bar{\frakq}\)에 따른 선형독립인 잉여류가 여섯 개이므로 노름은 \(P^{(2)}\)의 제곱이다. 따라서 여기서는 초등인자형의 \(p^g\)제곱이 나타나지 않는다.

첫 예는 비완전체를 계수영역으로 삼으면 대수적 수체에서와 똑같이 1보다 높은 차수의 소아이디얼들이 나타날 수 있고, 이때 \(p^g\)를 소아이디얼의 차수라고 불러야 함을 보여 준다. 뒤의 예들은 분기아이디얼의 유사물, 곧 한 소수가 한 소아이디얼의 더 높은 거듭제곱인 으뜸아이디얼로 나누어지는 현상으로 보아야 한다. 주 10)에서 보았듯 초등인자형은 대수적 수체에서 한 아이디얼로 나누어지는 가장 작은 유리정수에 대응하기 때문이다.

\subsection*{§ 7. 절대소아이디얼}

\(P\)의 계수를 갖는 한 소아이디얼이 \(P\)를 확대하여 얻을 수 있는 대수적으로 닫힌 체 \(A\)에서도 소아이디얼로 남으면 이를 \emph{절대소아이디얼}이라 한다.\footnote{잘 알려졌듯 대수적으로 닫힌 체란 한 부정원의 모든 다항식이 일차인자들로 분해되는 체이다. 임의의 체에 속하는 대수적으로 닫힌 체의 존재와 본질적 유일성은 Steinitz를 보라.} 이에 대응하여 \(P\)의 계수를 갖는 소함수 \(P\)가 \(A\)에서도 소함수로 남으면 이를 \emph{절대소함수}, 곧 절대기약다항식이라 한다. 절대소아이디얼은 다음 정리로 특징지어진다.

\textbf{정리 XV.} \emph{한 절대소아이디얼의 필요충분조건은 \(u\)들에 관하여 다항식이고 원시적이라고 가정할 수 있는 그 초등인자형이 \(x\)와 \(u\)의 다항식으로서 절대소함수인 것이다. 초등인자형은 노름과 같으므로 이 조건을 노름으로 진술해도 된다.}

증명을 위해 먼저 계수영역을 대수적으로 확대해도 일반적으로 노름과 초등인자형은 그대로 유지됨을 주목한다. 개별 인자 \(R^{(i)}\), 또는 \(E^{(i)}\)은 다항식 의미의 최대공약인자로 정의되며, 이 성질은 계수영역의 대수적 확대 아래에서도 유지되기 때문이다. 따라서 \(P\)에서 \(A\)로 넘어갈 때도 \(P^{(i)}\)은 \(\frakp\)의 초등인자형으로 남는다. 이는 계수영역을 완전체로 바꾸는 것이다. 대수적으로 닫힌 체는 정의에서 곧바로 알 수 있듯 완전체이기 때문이다. 따라서 정리 XIII에 따라 \(\frakp\)가 절대소아이디얼일 필요충분조건은 \(P^{(i)}\)이 \(A(u)\)에 관하여 소함수인 것이다. \(P^{(i)}\)을 \(u\)들에 관하여 다항식이고 원시적이라고 가정하면, 이는 \(x\)와 \(u\)의 다항식으로서 절대소함수라는 뜻이다. 동시에 \(A\)를 계수영역으로 삼는 순간 정리 XIII에 따라 \(P^{(i)}\)은 \(\frakp\)의 노름과도 같다. 앞서 말한 바에 따라 \(P\)를 바탕으로 삼아도 똑같이 성립한다. 이로써 정리 XV를 증명하였다.

이제 한 (유한) 대수적 수체 \(\mathfrak K\)의 계수를 갖는 절대소함수 \(P\)에 대해서는 다음 정리가 성립한다.\footnote{A. Ostrowski, 앞의 논문(주 21), 296쪽의 보조정리. 더 간단한 증명은 E. Noether, \emph{Ein algebraisches Kriterium für absolute Irreduzibilität}, Math. Ann. 85 (1922), 26--33쪽, 제7번에 있다.} \(\mathfrak K\)의 유한개 소아이디얼을 제외하면, \(\mathfrak K\)의 모든 소아이디얼을 법으로 해도 \(P\)는 소함수, 더 정확히는 절대소함수로 남는다. 이 정리를 절대소아이디얼로 옮기는 데에는 다음 정리가 쓰인다.

\textbf{정리 XVI.} \emph{계수영역으로 한 체 대신 한 (유한) 대수적 수체 \(\mathfrak K\)의 모든 대수적 정수들의 환 \(\mathfrak o^*\)을 삼으면, 한 다항식아이디얼의 노름과 초등인자형은 \(\mathfrak K\)의 유한개 소아이디얼을 제외한 모든 소아이디얼을 법으로 해도 노름과 초등인자형으로 그대로 유지된다.}

증명을 위해 § 1, 1의 다항식영역을 조금 바꾸어야 한다. \(\bar{\frakS}\)는 \(y_1,\ldots,y_n\)의 모든 다항식 가운데 계수가 \(\mathfrak o^*\), 곧 \(\mathfrak K\)의 대수적 정수인 것들로 이루어진다고 하자. \(\frakS\)의 계수영역은 \(\mathfrak o^*[u]\), 곧 계수가 \(\mathfrak o^*\)인 \(u\)들의 모든 다항식으로 이루어져 있다고 하자. \(\bar{\frakm}\)이 \(\bar{\frakS}\)의 아이디얼이면 변환 (1)에 따라 \(\frakS\)의 변환된 아이디얼 \(\frakm\)으로 간다. 노름을 만들 때 \(\mathfrak o^*[u]\)의 유한개 다항식만 분모에 나타남을 보이면 정리 XVI가 곧바로 따른다.

우선 § 1, 5의 모듈표현 (3)으로 개별노름을 만드는 것은 \(x_{i-1}\)에 관하여 정칙인 다항식
\(
C^{(i-1)}=U_{i-1}(u)x_{i-1}^{k}+\hbox{낮은 차수의 항들}
\)
을 법으로 \(x_{i-1}\)의 거듭제곱들을 환원하는 것으로 귀착됨을 주목한다.\footnote{H.-N., § 4를 보라.} 여기서 \(C^{(i-1)}\)의 모든 계수, 특히 \(U_{i-1}\)을 \(\mathfrak o^*[u]\)의 원소라고 가정할 수 있다. \(x_1,\ldots,x_{i-1}\)을 차례로 환원하므로, \(x\)에 관하여 다항식인 모듈표현은 분모에 나타나는 거듭제곱곱 \(U_1^{\lambda_1}\cdots U_{i-1}^{\lambda_{i-1}}\)만 제외하면 \(\mathfrak o^*[u]\)에서도 다항식이다. 그런데 \(u\)의 다항식으로 본 곱 \(U_1U_2\cdots U_n\)의 계수들은 \(\mathfrak K\)의 한 아이디얼 \(\mathfrak r^*\)을 정하고, 이는 \(\mathfrak K\)의 유한개 소아이디얼로만 나누어진다. 이 유한개와 다른 \(\frakp^*\)를 택하고 \(\frakp^*\)에 따른 잉여류체를 계수영역으로 삼으면, \(\mathfrak o^*\)의 각 수를 \(\frakp^*\)에 따른 그 잉여류로 바꾸기만 하여 원래 모듈표현이 유지된다.

또 노름과 초등인자형은 \(\mathfrak o^*[u]\)의 인자를 곱하는 차이까지만 정해지므로, \(\mathfrak o^*[u]\)에 관해서도 \(\frakm\)으로 나누어진다고 가정할 수 있다. 이 가정 아래 이를테면
\(
R^{(i)}=V_i(u)x_i^{l}+\hbox{낮은 차수의 항들}
\)
이라 하자. 그러면 \(C^{(i-1)}\)이 정하는 계수 \(\rho\)의 가군 \(\mathfrak M_{i-1}^*\)의 모든 \(\rho\)차 행렬식이 나누어지는 것을 다룰 때, 분모에는 \(U_1^{\lambda_1}\cdots U_{i-1}^{\lambda_{i-1}}\) 외에 \(V_i\)의 거듭제곱만 나타난다. 따라서 \(\mathfrak K\)에서 \(V_1V_2\cdots V_n\)이 정하는 아이디얼에 포함되는 유한개 소아이디얼과도 다른 \(\frakp^*\)를 택하면, \(\frakp^*\)에 따른 잉여류체를 계수영역으로 삼아도 \(R^{(i)}\)은 이 행렬식들의 공약인자로 남고 \(\mathfrak M_{i-1}^*\)의 계수도 유지된다. \(C^{(i)}\)이 \(\mathfrak M_{i-1}^*\)의 한 \(\rho\)차 행렬식이었기 때문이다.

끝으로 각 \(R^{(i)}\)이 최대공약인자로 남도록 \(\frakp^*\)를 더 골라야 한다. \(D^{(i)}\)가 \(\mathfrak M_{i-1}^*\)의 모든 \(\rho\)차 행렬식을 돌게 하자. 가정에 따라
\(
U_1^{\lambda_1}\cdots U_{i-1}^{\lambda_{i-1}}V_i^\sigma D^{(i)}
=R^{(i)}T^{(i)}
\)
이고, \(T^{(i)}\)들은 다항식 의미에서 \(x\)를 포함하는 공약인자를 갖지 않는다. 따라서 모든 \(T^{(i)}\)에서 유도된 아이디얼은
\(
G^{(i+1)}=W_i(u)x_{i+1}^{m}+\hbox{낮은 차수의 항들},
\qquad W_i\ne0
\)
인 다항식을 포함한다. \(G^{(i+1)}\)을 \(x_{i+1}\)에 관하여 정칙이라고 가정할 수 있음은 변환 (1)에 새 부정원들을 계수로 하는 \(x_{i+1},\ldots,x_n\)의 변환을 합성하면 알 수 있다. 그러므로 \(\mathfrak K\)에서 \(W_1W_2\cdots W_n\)이 정하는 아이디얼에 포함되는 유한개 소아이디얼과도 다른 \(\frakp^*\)를 택하면, \(\frakp^*\)에 따른 잉여류체를 계수영역으로 삼아도 \(R_\frakm\)은 실제로 노름으로 유지된다. 완전히 같은 방식으로 \(\frakp^*\)를 더 골라 \(E_\frakm\)도 초등인자형으로 유지되게 할 수 있다. 이로써 정리 XVI를 증명하였다.

정리 XV와 XVI, 그리고 절대소함수에 관한 위 정리에서 그 일반화인 다음 정리가 따른다.

\textbf{정리 XVII.} \emph{\(\frakp\)가 한 (유한) 대수적 수체 \(\mathfrak K\)의 대수적 정수들을 계수로 갖는 절대소아이디얼이면, \(\mathfrak K\)의 유한개 소아이디얼을 제외한 모든 소아이디얼을 법으로 해도 \(\frakp\)는 소아이디얼, 더 정확히는 절대소아이디얼로 남는다.}

증명을 위해 \(\frakp\)의 노름 \(R_\frakp\)를 정리 XVI에서처럼 \(\mathfrak o^*[u]\)에 관해서도 \(\frakp\)로 나누어지도록 택하자. 그러면
\(
R_\frakp=T(u)P^{(i)}(x)
\)
이고, \(P^{(i)}(x)\)는 \(u\)들에 관하여 다항식이고 원시적이라고 가정할 수 있다. 따라서 정리 XV에 따라 \(P^{(i)}\)은 \(\mathfrak K\)의 계수를 갖는 \(u\)와 \(x\)의 다항식으로 보았을 때 절대소함수이다. 절대소함수에 관한 정리에 따라, \(\mathfrak K\)의 유한개 소아이디얼을 제외하면 모든 소아이디얼을 법으로 해도 \(P^{(i)}\)은 이 성질을 유지한다. 따라서 이 유한개 소아이디얼과 다르고, 정리 XVI에 나오는 또 다른 유한개 소아이디얼과도 다른 \(\frakp^*\)를 택하자. 그러면 \(P^{(i)}\)의 \(\mathfrak o^*\)-계수들을 \(\frakp^*\)에 따른 대응 잉여류들로 바꾸었을 때, \(P^{(i)}\)은 \(\frakp^*\)에 따른 잉여류체를 계수영역 \(P\)로 삼은 \(\frakp\)의 노름이고, 이 노름은 절대소함수이다. 따라서 \(P\)를 계수영역으로 삼아도 정리 XV에 따라 \(\frakp\)는 절대소아이디얼로 남는다. 이로써 정리 XVII를 증명하였다.

\begin{center}
괴팅겐, 1923년 5월 8일.\\[1em]
(1923년 3월 9일 접수.)
\end{center}

\clearpage
